\documentclass[11pt]{article}
\usepackage[a4paper,margin=27mm]{geometry}
\usepackage{amsmath,amssymb,amsthm,mathtools,booktabs,array,longtable}
\usepackage[T1]{fontenc}
\usepackage{lmodern}
\usepackage[colorlinks=true,linkcolor=blue,citecolor=blue,urlcolor=blue,linktoc=all]{hyperref}
\usepackage{microtype}
\allowdisplaybreaks
\setlength{\emergencystretch}{2em}
\numberwithin{equation}{section}
\newtheorem{theorem}{Theorem}[section]
\newtheorem{lemma}[theorem]{Lemma}
\newtheorem{proposition}[theorem]{Proposition}
\newtheorem{corollary}[theorem]{Corollary}
\theoremstyle{definition}
\newtheorem{definition}[theorem]{Definition}
\newtheorem{construction}[theorem]{Construction}
\theoremstyle{remark}
\newtheorem{remark}[theorem]{Remark}
\DeclareMathOperator{\tr}{tr}
\DeclareMathOperator{\rank}{rank}
\DeclareMathOperator{\diag}{diag}
\DeclareMathOperator{\Rea}{Re}
\DeclareMathOperator{\Ima}{Im}
\title{Exact complex coherence in codimension two}
\author{Xuan Li}
\hypersetup{pdftitle={Exact complex coherence in codimension two},pdfauthor={Xuan Li}}
\date{}
\begin{document}
\maketitle
\begin{abstract}
We determine the minimum coherence of $N$ unit vectors in
$\mathbb C^{N-2}$ for every integer $N\ge180$. The optimum is given by
an explicit four-block construction with balanced integer multiplicities.
The same optimum holds among all complex $N\times N$ matrices with
diagonal entries one and rank at most $N-2$. Every extremizer is equivalent to the
balanced four-block Gram matrix. The proof uses kernel duality to reduce
the matrix problem to a geometric comparison in $\mathbb C^2$.
Any competing matrix whose coherence is no larger than that of the
construction admits a normalized two-dimensional left-kernel
representation with every column near one of four directions.
A strict local comparison then determines all equality cases.
Rational arithmetic certificates establish the finite inequalities needed
for the range $N\ge180$. For $4\nmid N$, the exact value improves the
classical codimension-two bound by a positive term of order $N^{-3}$.
\end{abstract}
\setcounter{tocdepth}{2}
\tableofcontents
\clearpage
\section{The optimization problem and the main result}
The coherence of a family of unit vectors is the largest modulus of an
inner product between distinct vectors. Minimizing coherence is a
packing problem for complex lines: it asks how well a prescribed number
of lines can be separated in a given dimension.
All inner products are conjugate-linear in the first variable:
$\langle v,w\rangle=v^*w$. For integers $1\leq n<N$ define
\[
 \mu_{\mathbb C}(n,N)=
 \min_{\substack{v_1,\ldots,v_N\in\mathbb C^n\\\|v_i\|=1}}
 \max_{i\ne j}|v_i^*v_j|.
\]
The minimum exists: the product of $N$ unit spheres is compact and the
maximum of the finitely many continuous inner-product moduli is continuous.

We study the case $n=N-2$. The associated Gram matrices have a kernel
of dimension at least two; a two-dimensional subspace of that kernel
provides the geometric representation used below. When $4\mid N$, the
classical attaining construction corresponds to equal replication of
four equiangular lines in this representation~\cite{BC,Glazyrin}.
For the remaining cardinalities, equal replication is impossible.
We construct a candidate using balanced integer multiplicities and
deform the four-line configuration accordingly. We prove that this
candidate is globally optimal among all vector families for every
$N\ge180$, and determine all equality cases.

The kernel argument also applies to matrices that need not initially
be Gram matrices. We therefore consider the two matrix problems
\[
 a_N=\min_{\substack{A\in\mathbb C^{N\times N}\\A_{ii}=1,\ \rank A\leq N-2}}
       \max_{i\ne j}|A_{ij}|,\qquad
 b_N=\inf_{\substack{A\in\mathbb C^{N\times N}\\A_{ii}=1,\ \rank A=N-2}}
       \max_{i\ne j}|A_{ij}|.
\]
For $N\geq3$ the first minimum also exists. The all-ones matrix is feasible
with objective one, so we can restrict to the compact set with all
off-diagonal moduli at most one. Diagonal and rank-at-most conditions are
closed. Moreover $a_N>0$, because a diagonal-one matrix with zero
off-diagonal entries has rank $N$. No compactness of the exact-rank set is
assumed; attainment of $b_N$ will follow from the construction.

Set $\alpha=(1+\sqrt3)/4$. A four-tuple
$\boldsymbol m=(m_1,m_2,m_3,m_4)$ of positive integers with sum $N$ is \emph{balanced} if its
entries differ by at most one. Such a tuple is unique up to permutation.
Put $p_a=m_a/N$ and denote by $\lambda_{\boldsymbol p}$ the root specified
in the following theorem.

\begin{theorem}\label{thm:main}
For every integer $N\geq180$ and every balanced four-tuple
$\boldsymbol m$ of positive integers with sum $N$, the equation
\begin{equation}\label{eq:main-root}
 2\lambda=\sum_{a=1}^4\frac{p_a(1-3p_a)}{\lambda-2p_a}
\end{equation}
has a unique root in $(2/3,\alpha]$, and
\begin{equation}\label{eq:main-value}
 \mu_{\mathbb C}(N-2,N)=a_N=b_N=
 \frac1{N\lambda_{\boldsymbol p}-1}.
\end{equation}
The exact-rank infimum is attained. Every minimizing matrix in the
rank-at-most problem is, up to index permutation, diagonal unitary
switching and entrywise conjugation, the balanced four-block Gram matrix
in Proposition~\ref{prop:construction}. In particular it is positive semidefinite,
has rank $N-2$, and all off-diagonal entries have the same modulus.
Vector realizations differ additionally by a common unitary isometry.
\end{theorem}

For $4\mid N$ the value is $1/(\alpha N-1)$, the known divisible case.
For $N\equiv2\pmod4$ the formula becomes
\begin{equation}\label{eq:residue-two}
 \mu_{\mathbb C}(N-2,N)=
 \frac4{N+\sqrt{3N^2-8}-4}.
\end{equation}

\begin{corollary}\label{cor:asymptotic}
For each fixed $r\in\{0,1,2,3\}$, as $N\longrightarrow\infty$ through
$N\equiv r\pmod4$,
\[
 \mu_{\mathbb C}(N-2,N)=\frac1{\alpha N-1}
 +\frac{r(4-r)}{4\sqrt3\,\alpha^2}\,N^{-3}+O(N^{-4}).
\]
For $r\ne0$ the correction is positive.
\end{corollary}
\begin{proof}
Write $p_a=1/4+\epsilon_a$, $\sum_a\epsilon_a=0$ and
$\sum_a\epsilon_a^2=r(4-r)/(4N^2)$. Define
$\Phi_{\boldsymbol p}(\lambda)=4[2\lambda-
\sum_ap_a(1-3p_a)/(\lambda-2p_a)]$ and put $h=\alpha-1/2$.
Polynomial division gives
\[
 \Phi_{\boldsymbol p}(\alpha)
 =4\sum_a\frac{\epsilon_a^2}{h-2\epsilon_a},\qquad
 \Phi'_{(1/4,\ldots,1/4)}(\alpha)=\frac{4\sqrt3}{h}.
\]
The derivative is positive near uniform weights; Taylor's formula gives
$\alpha-\lambda_{\boldsymbol p}
=\sum_a\epsilon_a^2/\sqrt3+O(N^{-3})$.
Substitution in \eqref{eq:main-value} proves the assertion.
\end{proof}

The correction in Corollary~\ref{cor:asymptotic} measures the cost of the
integer multiplicities. It vanishes when equal replication is possible
and is positive in each of the other residue classes.

The proof begins with a trace inequality relating the coherence of a
matrix to an energy associated with a two-dimensional subspace of its
left kernel. A matrix that matches or improves on the construction must
have kernel energy close to the classical upper bound. We use the
resulting shared deficit bounds to place every kernel
column in one of four clusters and to restrict their integer counts.
A local comparison then controls both motion of the cluster centers and
variation within each cluster, including the complex phases. Its strict
inequality forces an extremizer to coincide with the corresponding
four-block candidate. Integer balancing selects the optimal counts,
and equality in the trace inequality recovers the full matrix.

The explicit construction and matrix argument are given in
Sections~\ref{sec:construction} and~\ref{sec:proof}. The quantitative
localization and local comparison use rational inequalities on closed
domains, whose certificates are described in Section~\ref{app:certificates}.
The threshold $180$ is sufficient; its minimality remains open.

\section{Relation to earlier work}\label{sec:related}

Bukh and Cox~\cite[Theorems~2(b) and~3(b)]{BC} proved a bound for general
matrices. For a complex matrix $A$ of order
$N$ with diagonal one and rank at most $N-k$, it gives
\begin{equation}\label{eq:bc-literature}
 \max_{i\ne j}|A_{ij}|\geq\frac{1}{\alpha_kN-1},
 \qquad
 \alpha_k=\frac{1+(k-1)\sqrt{k+1}}{k^2}.
\end{equation}
Although their extremal quantity is defined using exact rank, the proof
of their Theorem~10 already permits rank at most $N-k$. Equality in
\eqref{eq:bc-literature} requires a maximal equiangular tight frame in
$\mathbb C^k$ and $k^2\mid N$; uniform replication supplies the matching
construction. In codimension two this is the tetrahedral configuration
of four complex lines. Consequently, the exact value
$((1+\sqrt3)N/4-1)^{-1}$ when $4\mid N$ is an established result.

Their corrigendum~\cite{BCErratum} specifies the complex equality case:
equal mass is assigned to each entire
unit-phase orbit, rather than to a pair of antipodal representatives.
This correction preserves the numerical bound and its divisibility
condition. We use the corrected projective formulation throughout.

Magsino, Mixon and Parshall~\cite[Lemma~1 and Theorem~4]{MMP} give a
Delsarte-style proof through an auxiliary polynomial and tight-frame
duality. Glazyrin~\cite[Theorem~7]{Glazyrin} develops a general matrix
duality for conjugate energies. His replicated-code construction and
equality analysis~\cite[Theorem~10 and the subsequent discussion]{Glazyrin}
also show that an extremizing general matrix is a Gram matrix, with
uniqueness inherited from the maximal projective simplex. In particular,
the divisible codimension-two case already has this stronger equality
description. We use this duality, the auxiliary-polynomial estimate,
phase alignment at equality, and the tetrahedral model in the proof below.

Two closely related optimization problems help explain the dual
functional. Cahill and Casazza~\cite[Proposition~2.1 and Theorems~2.2,
2.4]{CahillCasazza} study total coherence of Parseval frames, its
invariance under Naimark complementation, and restrictions on maximizing
frames. For columns $u_i$ of a matrix $U\in\mathbb C^{2\times N}$ with
$UU^*=I_2$, write
\[
 \operatorname{TC}(U)=\sum_{i\ne j}|u_i^*u_j|.
\]
The quasimaximal relative projection constant of Foucart and
Skrzypek~\cite[Section~2]{FoucartSkrzypek} is, in this notation,
\[
 q_{\mathbb C}(2,N)
 =\frac{\max_{UU^*=I_2}\operatorname{TC}(U)+2}{N}.
\]
Their maximal relative projection constant additionally optimizes
independent nonnegative weights $t_i$, with $\sum_i t_i^2=1$, in
$\sum_{i,j}t_it_j|u_i^*u_j|$. These distinctions matter when transferring
an exact value or an equality statement to the coherence problem.

Kobos~\cite[Theorem~4.2]{Kobos} proves a stability theorem: a sequence
approaching the maximal ETF bound admits, after
passing to a subsequence, clustering around maximal ETF directions in
a weighted mean-square sense; an exceptional set may have vanishing
total weight. His analysis also treats phases and residuals. To obtain
the finite conclusion here, we need bounds that place every column in
a controlled neighborhood at each cardinality under consideration.
The subsequent local comparison must eliminate all within-cluster
residuals. These are the quantitative steps developed in this paper.

Der\k{e}gowska and Lewandowska~\cite[Theorem~3.6]{DLWelch} determine
relative projection constants at replicated maximal-ETF cardinalities.
Kumar, Mohar, Mojallal and Pragada~\cite[Theorem~1.3]{KMMP} establish
eventual stabilization of maximal relative projection constants in the
real setting. The codimension-two theorem of Der\k{e}gowska and
Lewandowska~\cite[Theorem~2.4]{DLCodim} instead controls a real
Banach--Mazur distance through a small projection constant. These
results concern different objectives from the finite complex
off-diagonal minimax problem in Theorem~\ref{thm:main}.

Theorem~\ref{thm:main} extends the exact value and complete equality
classification to the nondivisible cardinalities in the range
$N\geq180$, within the full class of complex unit-diagonal matrices of
rank at most $N-2$. It identifies the balanced integer multiplicities,
the corresponding deformation of the four-line configuration, and all
extremizers. Quantitative local rigidity and computer-assisted existence
arguments also occur in projective coding theory, for example
in~\cite{CKM}. Here the local comparison is combined with control of
every kernel column to obtain a global minimax theorem uniformly over
the stated cardinalities.

\section{The attaining vectors and their stress}\label{sec:construction}
We give the construction for every \emph{admissible} profile
$\boldsymbol p=(p_1,\ldots,p_4)$, meaning $0<p_a<1/3$ and $\sum p_a=1$.
Let $\lambda=\lambda_{\boldsymbol p}$ be the root of
\eqref{eq:main-root}, and put
\begin{equation}\label{eq:construction-data}
 d_a=\lambda-p_a,\qquad w_a=p_a/d_a,\qquad
 X_{aa}=1,\quad X_{ab}=2w_aw_b-1\quad(a\ne b).
\end{equation}
The scalar equation has a unique root in $(2/3,\alpha]$, and $X$ is a
real positive semidefinite matrix of rank three. The algebraic proofs,
including strict integer balancing, are in Section~\ref{app:candidate}.
Realize $X$ by four unit Bloch vectors $x_a\in\mathbb R^3$ and choose
unit spinors $u_a\in\mathbb C^2$ with
$|u_a^*u_b|^2=(1+x_a\cdot x_b)/2$.

\begin{proposition}[Explicit vector factor]\label{prop:construction}
Let $m_a=Np_a$ be positive integers. Form the $2\times4$ matrix
$W=[\sqrt{d_1}u_1\ \cdots\ \sqrt{d_4}u_4]$ and put
$B=W^*W$, $A_p=\lambda I_4-B$. Let $J$ be the $N\times4$ block
isometry, with entry $1/\sqrt{m_a}$ in column $a$ on block $a$.
If the columns of $R\in\mathbb C^{N\times(N-4)}$ form an orthonormal
basis of $\operatorname{range}(J)^\perp$, then
\begin{equation}\label{eq:explicit-vectors}
 V_p=\frac{1}{\sqrt{N\lambda-1}}
 \begin{pmatrix}\sqrt{N\lambda}\,R^*\\ \sqrt N\,WJ^*\end{pmatrix}
 \in\mathbb C^{(N-2)\times N}
\end{equation}
has unit columns, rank $N-2$, and coherence $(N\lambda-1)^{-1}$.
Its Gram matrix is
\begin{equation}\label{eq:attaining-gram}
 S_p=NJA_pJ^*,\qquad
 G_p=V_p^*V_p=\frac{N\lambda I_N-S_p}{N\lambda-1}.
\end{equation}
For a fixed ordered profile, the Gram matrix $G_p$ is unique up to
diagonal unitary conjugation and entrywise conjugation; permuting the
profile permutes its blocks.
\end{proposition}
\begin{proof}
The Bloch identity gives $B_{aa}=\lambda-p_a$ and
$|B_{ab}|=\sqrt{p_ap_b}$ for $a\ne b$. The matrix $B$ is positive
semidefinite of rank two. Thus $A_p$ has top eigenvalue $\lambda$
with multiplicity exactly two. Every entry of $S_p$ has modulus one,
and its diagonal entries equal one. Consequently $N\lambda I-S_p$
is positive semidefinite with exactly two zero eigenvalues and constant
diagonal $N\lambda-1$. Moreover
\[
 N\lambda I-S_p=N\lambda(I-JJ^*)+NJBJ^*.
\]
Using $RR^*=I-JJ^*$ proves \eqref{eq:explicit-vectors} and all the
claimed rank and modulus assertions. Orthogonal realizations of the
rank-three real Gram matrix $X$ differ by a rotation, or a rotation
followed by a reflection. On spinors these give a common unitary, or a
common unitary followed by conjugation; individual lift phases give
diagonal unitary conjugation. The choice of $R$ changes only the common
unitary factor of the vector realization.
\end{proof}

To compare the construction with arbitrary matrices, we associate a
positive $2\times2$ matrix $K$ to its kernel. We call $K$ the candidate
stress. The trace bound in Section~\ref{sec:proof} will convert an upper
bound for its kernel functional $\Gamma_K$ into a coherence lower bound.
Set
\[
 q_a=\frac{(1-\lambda-p_a)(\lambda-p_a)}
 {2(1-\lambda)(\lambda-2p_a)},\quad
 \ell_a=\frac{2q_a}{\lambda-p_a},\quad
 L=\diag(\ell_a),\quad T=L-LBL.
\]
Choose $U\in\mathbb C^{4\times2}$ with $U^*U=I$ and $BU=0$.
Writing $Y=[y_1\ \cdots\ y_N]\in\mathbb C^{2\times N}$, define
\begin{equation}\label{eq:candidate-stress}
 Z=\sqrt{N/2}\,U^*J^*,\qquad
 K=\frac{2U^*TU}{\tr T},\qquad
 \Gamma_K(Y)=\sum_{i,j}|y_i^*Ky_j|.
\end{equation}
The identities $q_a>0$, $\sum q_a=1$, $Xq=0$ imply $WLW^*=I$.
Hence $T\succeq0$ has rank two and range $\ker B$, so $K\succ0$ and
$\tr K=2$. Section~\ref{app:candidate} verifies that
\begin{equation}\label{eq:stress-identities}
 ZZ^*=\frac N2I,\quad ZS_p=N\lambda Z,\quad
 \operatorname{phase}(Z^*KZ)=S_p,\quad
 \Gamma_K(Z)=N^2\lambda.
\end{equation}
In particular, $ZG_p=0$, so $Z$ is the normalized left-kernel frame of
the candidate Gram matrix, with columns constant on each block.
All overlaps in the phase identity are nonzero, including those within
a block. Holding $K$ fixed, $Z$ is stationary for
$\Gamma_K(Y)-N\lambda\tr(KYY^*)$.

Finally, strict convexity of $p\mapsto(\lambda-2p)^{-1}$ in the scalar
equation shows that transferring one count from a larger block to a
block smaller by at least two strictly increases $\lambda_p$.
Thus
\begin{equation}\label{eq:strict-balancing}
 \lambda_p\le\lambda_{\rm bal},
 \qquad\text{with equality exactly for balanced integer counts.}
\end{equation}
After localization, this comparison selects the balanced profile from
the possible cluster counts.

\section{From arbitrary matrices to the unique equality class}\label{sec:proof}
The left kernel of a matrix in our rank-constrained class supplies the
comparison functional. A positive stress gives a trace bound whose
equality condition determines the matrix entries.

\begin{lemma}[Left-kernel trace bound]\label{lem:trace}
Let $A\in\mathbb C^{N\times N}$ have diagonal one, rank at most $N-2$,
and coherence $\mu=\max_{i\ne j}|A_{ij}|$. Choose
$Y\in\mathbb C^{2\times N}$ with rows in the left kernel of $A$ and
normalization $YY^*=NI_2/2$, so that $YA=0$.
For every $K\succ0$ with $\tr K=2$,
\begin{equation}\label{eq:trace-bound}
 N(1+1/\mu)\le\Gamma_K(Y).
\end{equation}
If equality holds and all off-diagonal entries of $H=Y^*KY$ are nonzero,
then $A_{ij}=-\mu H_{ij}/|H_{ij}|$ for every $i\ne j$.
\end{lemma}
\begin{proof}
Take two orthonormal columns $Q$ in $\ker A^*$ and put
$Y=\sqrt{N/2}\,Q^*$. Then $HA=0$, $H\succeq0$, and $\tr H=N$.
Also $\mu>0$, since the identity has rank $N$. Expanding
$\tr(HA)=0$, separating the diagonal terms, and taking real parts gives
\[
 N=-\sum_{i\ne j}\Rea(H_{ij}A_{ji})
 \le\mu\sum_{i\ne j}|H_{ij}|=\mu(\Gamma_K(Y)-N).
\]
If equality holds, every ordered term with $H_{ij}\ne0$ attains both
its modulus and real-part bounds. Thus
$A_{ji}=-\mu\overline{H_{ij}}/|H_{ij}|$, which is the stated formula
after exchanging the indices.
\end{proof}

\subsection{The quantitative comparison used in the proof}
For $Y$ as in Lemma~\ref{lem:trace}, write $y_i=r_i u_i$ and let $x_i\in S^2$
be the Bloch vector of $u_i$. If $r_i=0$, choose $u_i$ arbitrarily at
this stage. Put
\[
 a=\frac{2\sqrt3-1}{8},\quad k=\frac{6\sqrt3-9}{4},\quad L=\sum r_i,
 \qquad g(t)=(1-\sqrt t)(3/4+k\sqrt t)(\sqrt t-1/\sqrt3)^2.
\]
The radial, moment and angular deficits are
\begin{align}\label{eq:deficits}
 u&=a(N^2-L^2),&
 v&=\frac k4\left\|\sum_i r_i x_ix_i^T-\frac L3I\right\|_F^2,&
 w&=\sum_{i,j}r_ir_jg((1+x_i\cdot x_j)/2).
\end{align}
They are nonnegative and satisfy
$u+v+w\le\alpha N^2-\Gamma_I(Y)$, by the exact polynomial identity
in Section~\ref{app:global}. Applying Lemma~\ref{lem:trace} with $K=I$
to a competitor with $\mu\le U_N=(N\lambda_{\rm bal}-1)^{-1}$ gives
\begin{equation}\label{eq:budget}
 u+v+w\le M,
 \qquad M=\begin{cases}1/2,&N\text{ odd},\\11/19,&N\equiv2\pmod4.
 \end{cases}
\end{equation}
The scalar candidate estimates establishing these budgets are included
in Section~\ref{app:candidate}.

\begin{proposition}[Localization of every column]\label{prop:localization}
Suppose $N\ge180$, $4\nmid N$, $YY^*=NI_2/2$, and
$\Gamma_I(Y)\ge N^2\lambda_{\rm bal}$. There is a partition of all
$N$ columns into four nonempty groups with positive counts $m_a$.
Set $\delta_a=m_a-N/4$ and $C=\sum\delta_a^2$. Up to permutation and
simultaneous sign reversal, its possible profiles are
\begin{center}
\begin{tabular}{ccc}
\toprule
Residue & $C$ & $(\delta_1,\delta_2,\delta_3,\delta_4)$\\
\midrule
$N$ odd & $3/4$ & $(-1,-1,-1,3)/4$\\
$N$ odd & $11/4$ & $(-5,-1,3,3)/4$\\
$N\equiv2\pmod4$ & $1$ & $(-1,-1,1,1)/2$\\
$N\equiv2\pmod4$ & $3$ & $(-3,1,1,1)/2$\\
\bottomrule
\end{tabular}
\end{center}
On this same partition let $\bar x_a$ be the group means and
$S=\sum_a\sum_{i\in a}\|x_i-\bar x_a\|^2$.
Then $S\le31E/N$, where $E=\sum_{i,j}g((1+x_i\cdot x_j)/2)$.
The simultaneous total and pointwise estimates
\eqref{app:global:LU-LP}--\eqref{app:global:physical-inputs}
hold in one common preliminary alignment to the repeated tetrahedral
frame $F=(f_{a(i)})_i$, where $a(i)$ is the group containing column $i$
and the $f_a$ are unit tetrahedral spinors.
\end{proposition}
\begin{proof}
Section~\ref{app:global} keeps the shared budget $u+v+w\le M$
through anchor selection and partitions every column. Its centered count
estimate gives $C<5$ in residue two and $C<19/4$ in the odd residues,
which imply exactly the listed profiles. The forward refinement on
this same partition gives the scatter and simultaneous polar estimates.
The finite inequalities used in these steps are specified in
Section~\ref{app:certificates}.
\end{proof}

The scales explain the role of the local step. The shared deficit is
$O(1)$, uniformly in $N$, and the count variance $C$ is bounded.
Localization and the candidate estimates place both $Y$ and the
candidate $Z$ with the same profile within $O(N^{-1/2})$ of $F$ in the
Frobenius norm. The triangle inequality and minimization over the allowed
alignments therefore give $d=\|Y-Z\|_F=O(N^{-1/2})$ in the closest gauge.
The separate transport of pointwise bounds to that gauge is proved in
Section~\ref{app:local:angular}. The relative normal component is at most
$d/\sqrt{2N}=O(N^{-1})$. The horizontal reference
curvature has a negative gap bounded away from zero; candidate
perturbations and the higher-order loss tend to zero. This explains why
the local comparison becomes strict for large $N$. To obtain the
explicit threshold $180$, we bound both terms on the localized regions
and certify the resulting inequalities by rational arithmetic.

\begin{proposition}[Conditional local comparison]
\label{prop:conditional-local}
Fix one of the localized profiles, its candidate $Z$ and its positive
stress $K$. Put a competitor $Y$ with $YY^*=NI_2/2$ in the
closest unitary and individual-phase gauge to $Z$, allowing conjugation
of the candidate, and set $d=\|Y-Z\|_F$. Suppose all validity conditions
in Section~\ref{app:local}, including the denominator, positivity, phase,
overlap and closest-gauge transport conditions, hold throughout a closed
region containing this competitor. Let $c_m,c_r$ be valid lower bounds
for the curvature gaps in the mean and residual directions, and $J,L$
valid upper bounds for
the mixed and pure normal terms, measured in the original Frobenius
metric as in \eqref{eq:main-quadratic}. Let
$A_{\rm rem},B_{\rm rem}\ge0$ supply one bound for the full Taylor
remainder using the common tangent fraction defined below. Define
$g_{\rm normal}$ as the negative of the maximum of the quadratic upper
bound in \eqref{eq:main-quadratic} over the normalized components, with
the normal fraction restricted by the bound for this region. If
\[
 g_{\rm normal}-\frac{A_{\rm rem}+
          \sqrt{A_{\rm rem}^2+B_{\rm rem}^2}}2\ge\varepsilon>0,
\]
then
\begin{equation}\label{eq:conditional-comparison}
 \Gamma_K(Y)\le N^2\lambda_p-\varepsilon N\|Y-Z\|_F^2.
\end{equation}
\end{proposition}
\begin{proof}
The case $d=0$ is immediate. For $d>0$, define the Lagrangian
$\mathcal L(Y)=\Gamma_K(Y)-N\lambda_p\tr(KYY^*)$ and its half-Hessian
$Q(H)=\tfrac12D^2\mathcal L(Z)[H,H]$, with $K$ fixed and derivatives
taken over the underlying real vector space. Put $\Delta=Y-Z$ and
$R_n=(\Delta\Delta^*/N)Z$. The horizontal component
$\Xi=\Delta+R_n$ splits orthogonally into cluster means $H_{\rm m}$ and
zero-sum residuals $H_{\rm r}$. Set $h=\|H_{\rm m}\|_F/d$,
$r=\|H_{\rm r}\|_F/d$, and $n=\|R_n\|_F/d$. Under the stated validity
conditions, the half-Hessian estimate is
\begin{equation}\label{eq:main-quadratic}
 \frac{Q(\Delta)}{Nd^2}\le-c_mh^2-c_rr^2+2Jnh+Ln^2,
 \quad h^2+r^2+n^2=1,\quad n\le d/\sqrt{2N}.
\end{equation}
The exact maximum of this quadratic upper bound over the allowed normal
fraction, including both branches, is $-g_{\rm normal}$; see
\eqref{app:local:normal-max}.

The remainder uses the distinct physical decomposition
$y_i=(1+t_i)z_i+\xi_i$, where $t_i\in\mathbb R$ and $z_i^*\xi_i=0$.
For the common tangent fraction
$\theta=(\sum_i\|\xi_i\|^2)^{1/2}/d\in[0,1]$, its complete upper bound is
$Nd^2[A_{\rm rem}\theta^2+B_{\rm rem}\theta\sqrt{1-\theta^2}]$.
Lemma~\ref{app:local:loss} bounds this by
$Nd^2(A_{\rm rem}+\sqrt{A_{\rm rem}^2+B_{\rm rem}^2})/2$.
The first variation of $\mathcal L$ vanishes at $Z$, and the constraint
term $N\lambda_p\tr(KYY^*)$ has the same value at $Y$ and $Z$ because
their covariances agree. Adding the quadratic and remainder bounds proves
\eqref{eq:conditional-comparison}. Section~\ref{app:local} derives the
bounds and their validity conditions; the next proposition certifies
those conditions on the regions supplied by localization.
\end{proof}

\begin{proposition}[Finite verification and propagation]
\label{prop:finite-acceptance}
The exact verification statement in Section~\ref{app:certificates:acceptance}
covers all closed budget cells of the four profiles. The local endpoints
are $180$ for $C\in\{3/4,11/4,3\}$ and $182$ for $C=1$.
Each cell is either excluded by a strict counting inequality or lies in a
closed region where all validity conditions of
Proposition~\ref{prop:conditional-local} hold and the quadratic gap exceeds
the remainder bound by more than $1/10000$.
For each fixed signed deviation vector $(\delta_a)$, these inequalities
remain valid for every larger integer $N$ for which
$m_a=N/4+\delta_a$ are integers.
\end{proposition}
\begin{proof}
The rational records and the conditions checked by the verifier are given in
Section~\ref{app:certificates:acceptance}.
Lemma~\ref{app:local:propagation} extends the verified inequalities by
fixing valid endpoint coefficients and bounding the reciprocal parameter
over the full subsequent interval.
\end{proof}

\begin{corollary}[Localization and strict comparison]
\label{prop:comparison}
Under the hypotheses of Proposition~\ref{prop:localization}, use its
partition and the closest gauge to the candidate with the same profile.
Then
\begin{equation}\label{eq:strict-comparison}
 \Gamma_K(Y)\le N^2\lambda_p-
                  \frac{N}{10000}\|Y-Z\|_F^2.
\end{equation}
\end{corollary}
\begin{proof}
Combine Propositions~\ref{prop:localization},
\ref{prop:conditional-local} and \ref{prop:finite-acceptance}.
Odd integers in the range start at $181$; the first actual residue-two
integer is $182$. Thus the two different local endpoints leave no
integer gap for nondivisible $N\ge180$.
\end{proof}

\subsection{Completion of the proof and equality cases}
We now prove Theorem~\ref{thm:main}. For $4\nmid N$, take any complex
$N\times N$ matrix $A$ with diagonal entries one, rank at most $N-2$,
and coherence $\mu\le U_N$. Choose $Y$ as in Lemma~\ref{lem:trace}. The closest
gauge is legitimate for a general matrix: replacing $Y$ by $UYD$ is
accompanied by $A\mapsto D^*AD$. It preserves $YA=0$, covariance,
diagonal, rank and coherence. Select the stress $K$ for this profile \emph{after}
this alignment. Lemma~\ref{lem:trace}, Corollary~\ref{prop:comparison}
and strict balancing give
\[
\begin{aligned}
 N(1+1/\mu)&\le\Gamma_K(Y)\\
 &\le N^2\lambda_p-\frac N{10000}\|Y-Z\|_F^2\\
 &\le N^2\lambda_{\rm bal}\le N(1+1/\mu).
\end{aligned}
\]
Every inequality is an equality. Therefore $Y=Z$ in the chosen gauge,
the counts are balanced, and $\mu=U_N$. All stress overlaps are nonzero,
so the ordered equality condition in Lemma~\ref{lem:trace} forces
\[
 A=(1+\mu)I-\mu S_p=G_p.
\]
Undoing the allowed transformations gives exactly the equality class
of the theorem.

For $4\mid N$, the classical scalar bound is sharp. For completeness,
the deficit identity at zero forces $r_i=1$, vanishing moment defect,
and $g((1+x_i\cdot x_j)/2)=0$ for every pair. Thus distinct Bloch
directions have inner product $-1/3$. At most four such directions exist.
Since $r_i=1$, tightness gives $\sum x_i=0$; together with the zero
moment defect this gives $\sum l_il_i^T=NI_4/4$, where
$l_i=(1,\sqrt3 x_i)^T/2$. The lifted directions
therefore span $\mathbb R^4$, so there are exactly four. They form an
orthonormal basis of $\mathbb R^4$, and tightness forces equal counts
$N/4$. Hence the kernel consists of equally repeated tetrahedral spinor
lines. Equality in the ordered trace bound with $K=I$ again fixes
$A=G_p$. This also proves the full matrix equality statement in the
divisible case.

Proposition~\ref{prop:construction} supplies an attaining rank-exactly
$N-2$ Gram matrix. Since the rank-at-most problem is a relaxation of
both other problems, all three optimal values coincide. Finally, two
full-rank vector factors with the same Gram matrix differ by a common
unitary on $\mathbb C^{N-2}$, completing the vector classification.

\section{Further questions and verification}
The theorem resolves the nondivisible cardinalities in an explicit
sufficient range. Determining the least threshold for the balanced
four-block construction remains open. In higher codimension, a related
question is whether the extremizers at nondivisible cardinalities arise
by deforming replicated maximal equiangular tight frames, when such
frames exist. An extension would require both a localization theorem
for the corresponding projective configuration and a strict comparison
for its possible integer multiplicities.

Two features of the present argument may be useful for these questions.
The common deficit budget couples radial and angular deviations during
localization. With the candidate stress fixed, the kernel functional
$\Gamma_K$ is exactly quadratic under positive radial rescalings of
the columns. The radial--tangential decomposition in the local
comparison uses this property to control the remainder. These features
reduce the loss in passing from a global bound to a finite neighborhood
of the candidate.

The proof uses computer assistance to verify the finite inequalities in
Sections~\ref{app:candidate}--\ref{app:certificates}. The supplementary
package provides rational certificates, an arithmetic verifier that checks
them without repeating the certificate search, and instructions for
reproducing the verification.
Section~\ref{app:certificates} specifies the verifier's scope and the
analytic arguments that extend the finite checks to the stated range.

\section{Algebra of the four-weight construction}\label{app:candidate}
\label{alg:section}

Throughout, $\langle u,v\rangle=u^*v$, and a collection of $N$ vectors is
written as a matrix with $N$ columns. Put $\alpha=(1+\sqrt3)/4$.
A profile is admissible if $0<p_a<1/3$ for $1\leq a\leq4$ and
$\sum_a p_a=1$. Unsubscripted matrix norms are operator norms.
We distinguish the kernel projector $\Pi$, the mass matrix
$D_{\mathrm{mass}}$, the Bloch Gram difference $\Delta_{\mathrm B}$,
and the scalar candidate point bound $C_{\mathrm{pt}}$.
The local Hessian error bound in Section~\ref{app:local} is denoted
$E_{\mathrm{Hess}}$.

\begin{lemma}[The scalar parameter]\label{alg:root}
For every admissible profile there is exactly one
$\lambda_p\in(2/3,\alpha]$ satisfying
\begin{equation}\label{alg:root-equation}
 2\lambda_p=\sum_a\frac{p_a(1-3p_a)}{\lambda_p-2p_a}.
\end{equation}
Moreover $\lambda_p<1-\max_a p_a$, and $\lambda_p=\alpha$ if and only if
all four weights are $1/4$.
\end{lemma}
\begin{proof}
For $\lambda\geq2/3$ all denominators below are positive. Define
\begin{align}
 \Phi_p(\lambda)
 &=\lambda(3\lambda-2)\sum_a\frac1{\lambda-2p_a}
                  -2(2\lambda-1)\label{alg:Phi}\\
 &=4\left[2\lambda-\sum_a\frac{p_a(1-3p_a)}{\lambda-2p_a}\right].
 \nonumber
\end{align}
Polynomial division and $\sum p_a=1$ give the second identity. Thus
\[
 \Phi_p(2/3)=-2/3,\qquad
 \Phi_p'(\lambda)=4\left[2+
 \sum_a\frac{p_a(1-3p_a)}{(\lambda-2p_a)^2}\right]>8.
\]
Strict convexity gives $\sum_a(\alpha-2p_a)^{-1}\geq4/(\alpha-1/2)$.
Since $\alpha(3\alpha-2)>0$, substitution yields
$\Phi_p(\alpha)\geq0$, with equality precisely at the uniform profile.
Continuity and strict monotonicity prove existence and uniqueness.
For the extra bound put $p=\max p_a$ and $L=1-p>2/3$. The other three
weights sum to $L$, so their denominators $L-2p_a$ sum to $L$. Therefore
\[
 \sum_a\frac1{L-2p_a}\geq\frac1{3L-2}+\frac9L,\qquad
 \Phi_p(L)\geq L(3L-2)\left[\frac1{3L-2}+\frac9L\right]
                 -2(2L-1)=8(3L-2)>0.
\]
Strict monotonicity implies $\lambda_p<L$.
\end{proof}


Abbreviate $\lambda=\lambda_p$, and put
\begin{equation}\label{alg:weights-X}
 d_a=\lambda-p_a,\quad w_a=p_a/d_a,\quad
 X_{aa}=1,\quad X_{ab}=2w_aw_b-1\quad(a\ne b).
\end{equation}
Here $X$ is an auxiliary four-point real Gram matrix, distinct from
the Bloch points of an arbitrary competitor.

\begin{lemma}[The auxiliary Gram matrix and its positive kernel]
\label{alg:bloch}
The matrix $X$ is positive semidefinite of rank three. Its kernel is
spanned by the strictly positive vector
\begin{equation}\label{alg:q}
 q_a=\frac{(1-\lambda-p_a)(\lambda-p_a)}
           {2(1-\lambda)(\lambda-2p_a)}.
\end{equation}
It has $\sum q_a=1$, and $q_a<1/2$ for every $a$.
\end{lemma}
\begin{proof}
We have $0<w_a<1$. Set
$D_X=2\operatorname{diag}(1-w_a^2)$, so
$X=D_X+2ww^T-\mathbf1\mathbf1^T$. Write
$A_s=\sum(1-w_a^2)^{-1}$, $B_s=\sum w_a/(1-w_a^2)$,
and $R=\sum(\lambda-2p_a)^{-1}$. Expanding the $2\times2$ determinant
for this rank-two update gives
\[
 \frac{\det X}{\det D_X}=(A_s-3)(1-A_s/2)+B_s^2/2.
\]
The identities
\[
 A_s-B_s=4-1/\lambda,\qquad A_s+B_s=2+\lambda R/2
\]
make $\det X=0$ equivalent to
$(3\lambda-2)(2+\lambda R/2)=8\lambda-5$, which is
$\Phi_p(\lambda)=0$. The matrix $D_X+2ww^T$ is positive definite.
On $\mathbf1^\perp$ its quadratic form remains unchanged after
subtracting $\mathbf1\mathbf1^T$. Thus $X$ has at least three strictly
positive eigenvalues; its remaining eigenvalue is zero.

Every factor in \eqref{alg:q} is positive by Lemma~\ref{alg:root}.
The normalization and kernel identities can be checked without an
eigenvalue computation. For arbitrary $\lambda$ in the displayed
domain, put $\beta=(2\lambda-1)/(2(1-\lambda))$. Division in
$\lambda-2p_a$ gives
\begin{equation}\label{alg:q-residual}
 \sum_aq_a-1=\sum_aw_aq_a-\beta
       =-\frac{\Phi_p(\lambda)}{8(1-\lambda)}.
\end{equation}
Also $2(1-w_a^2)q_a=1-2\beta w_a$. Consequently
$(Xq)_a=(1-2w_a)\Phi_p(\lambda)/(8(1-\lambda))$.
At the root, $\sum q_a=1$ and $Xq=0$.
Realize $X$ by unit vectors $x_a\in\mathbb R^3$. Then $\sum q_ax_a=0$,
so $q_a=\|\sum_{b\ne a}q_bx_b\|\leq1-q_a$.
Equality would force the three other points to coincide at $-x_a$,
contradicting rank three. Hence $q_a<1/2$.
\end{proof}

We use the elementary Bloch representation
$uu^*=(I+x\cdot\sigma)/2$ of a unit spinor $u\in\mathbb C^2$, where
\[
 \sigma_1=\begin{pmatrix}0&1\\1&0\end{pmatrix},\quad
 \sigma_2=\begin{pmatrix}0&-i\\i&0\end{pmatrix},\quad
 \sigma_3=\begin{pmatrix}1&0\\0&-1\end{pmatrix}.
\]
Multiplication and traces of these matrices give
\begin{equation}\label{alg:bloch-identity}
 |u^*v|^2=(1+x\cdot y)/2.
\end{equation}
Conversely a spinor for $x\in S^2$ is
\[
 u(x)=\left(\sqrt{(1+x_3)/2},
       \frac{x_1+i x_2}{\sqrt{2(1+x_3)}}\right)^T
\]
when $x_3\ne-1$, and $(0,1)^T$ at the south pole.

\begin{proposition}[An explicit attaining configuration]\label{alg:construction}
Let $m_a$ be positive integers with total $N$ and admissible
$p_a=m_a/N$. There are $N$ unit vectors in $\mathbb C^{N-2}$ whose Gram
matrix has rank exactly $N-2$ and constant off-diagonal absolute value
\begin{equation}\label{alg:mu-profile}
 \mu_p=\frac1{N\lambda_p-1}.
\end{equation}
To construct them, realize $X$ by real unit vectors $x_a$, choose
corresponding unit spinors $u_a$, and let $W\in\mathbb C^{2\times4}$
have columns $\sqrt{d_a}u_a$. Set
\begin{equation}\label{alg:core}
\begin{aligned}
 B&=W^*W,\qquad A_p=\lambda I_4-B,\\
 D_{\mathrm{mass}}&=\operatorname{diag}(\sqrt{p_a}),\qquad
 C=D_{\mathrm{mass}}^{-1}A_pD_{\mathrm{mass}}^{-1}.
\end{aligned}
\end{equation}
Let $J\in\mathbb R^{N\times4}$ be the isometric block lift, with
$J_{ia}=1/\sqrt{m_a}$ for $i$ in block $a$, and zero otherwise. Then
\begin{equation}\label{alg:signature-Gram}
 S=NJA_pJ^*,\qquad G_p=\frac{N\lambda I_N-S}{N\lambda-1}
\end{equation}
is the required Gram matrix. In particular $S_{ij}=C_{ab}$ for
$i\in a$, $j\in b$.
\end{proposition}
\begin{proof}
The Bloch identity gives $B_{aa}=\lambda-p_a$ and
$|B_{ab}|=\sqrt{d_ad_bw_aw_b}=\sqrt{p_ap_b}$ for $a\ne b$.
The matrix $B$ is positive semidefinite of rank two. Indeed rank one
would force all spinor lines, and hence all Bloch points, to coincide.
Thus $A_p$ has top eigenvalue $\lambda$ of multiplicity exactly two,
while $C_{aa}=1$ and $|C_{ab}|=1$. The spectrum of $S$ is the spectrum
of $NA_p$ together with $N-4$ zeros. Hence $N\lambda I-S$ is positive
semidefinite with exactly two zero eigenvalues. Its diagonal is
$N\lambda-1$, and every off-diagonal entry has modulus one.

An explicit vector factor is obtained as follows.
Choose $R\in\mathbb C^{N\times(N-4)}$ with orthonormal columns
spanning $\operatorname{range}(J)^\perp$. Then
\begin{equation}\label{alg:vector-factor}
 V_p=\frac1{\sqrt{N\lambda-1}}
       \begin{pmatrix}\sqrt{N\lambda}\,R^*\\ \sqrt N\,WJ^*\end{pmatrix}
       \in\mathbb C^{(N-2)\times N}
\end{equation}
satisfies $V_p^*V_p=G_p$, because
$RR^*=I-JJ^*$ and
$N\lambda I-S=N\lambda(I-JJ^*)+NJBJ^*$.
For $N=4$ the first block of rows is empty.
\end{proof}

\begin{remark}[Choices and equivalences]\label{alg:construction-choices}
The construction uses only the specified root, matrix square roots,
and the displayed spinor lift. The first three Bloch vectors can be
obtained by factoring their positive definite principal Gram matrix;
the fourth is then determined by its inner products. Positivity of
that principal matrix follows because the kernel of $X$ has no vector
with a zero coordinate except zero. Any two realizations of $X$ differ
by an orthogonal transformation. A proper rotation corresponds to a
common unitary on the spinors, and an improper rotation to a common
unitary followed by conjugation. One checks the rotation assertion
using $\cos(\theta/2)I-i\sin(\theta/2)n\cdot\sigma$; complex conjugation
reverses the second Bloch coordinate. Independent spinor phases give
diagonal unitary conjugations of $B,C,S,G_p$. Thus for an ordered
profile $G_p$ is unique up to diagonal unitary conjugation and possibly
entrywise conjugation. Permuting the profile permutes the blocks.
\end{remark}

\begin{proposition}[A positive stress with strict phase alignment]
\label{alg:stress}
For the above construction let
\begin{equation}\label{alg:stress-def}
 \ell_a=\frac{2q_a}{d_a}
 =\frac{1-\lambda-p_a}{(1-\lambda)(\lambda-2p_a)},\quad
 L=\operatorname{diag}(\ell_a),\quad T=L-LBL.
\end{equation}
Then $T\succeq0$ has rank two and range $\ker B$.
Choose $U\in\mathbb C^{4\times2}$ with $U^*U=I_2$ and $BU=0$, and set
\begin{equation}\label{alg:Z-K}
 z_a=\frac{U^*e_a}{\sqrt{2p_a}},\qquad
 Z=\sqrt{N/2}\,U^*J^*,\qquad K=\frac{2U^*TU}{\operatorname{tr}T}.
\end{equation}
The matrix $Z$ repeats $z_a$ exactly $m_a$ times and obeys
$ZZ^*=NI_2/2$ and $ZS=N\lambda Z$. Moreover $K\succ0$,
$\operatorname{tr}K=2$, and all entries of $Z^*KZ$ are nonzero with
entrywise phase $S$. Consequently
\begin{equation}\label{alg:candidate-value}
 \Gamma_K(Z):=\sum_{i,j}|z_i^*Kz_j|=N^2\lambda.
\end{equation}
With $K$ held fixed, $Z$ is a critical point of
\begin{equation}\label{alg:Lagrangian}
 \mathcal L(Y)=\Gamma_K(Y)-N\lambda\operatorname{tr}(KYY^*).
\end{equation}
\end{proposition}
\begin{proof}
Since $\sum q_ax_a=0$ and $\sum q_a=1$, the Bloch identity gives
$\sum q_au_au_a^*=I_2/2$, and hence $WLW^*=I_2$.
If $E=WL^{1/2}$, then $EE^*=I_2$, so $E^*E$ is a rank-two orthogonal
projector. Therefore
\[
 T=L^{1/2}(I-E^*E)L^{1/2}\succeq0,\qquad\operatorname{rank}T=2.
\]
Also $WT=WL-(WLW^*)WL=0$. Thus the range of $T$ is $\ker B$,
$U^*TU$ is positive definite, and $\operatorname{tr}(U^*TU)=\operatorname{tr}T$.
The exact entry formulas are
\[
 T_{ab}=\ell_a\ell_b\sqrt{p_ap_b}\,C_{ab}\quad(a\ne b),\qquad
 T_{aa}=\ell_a(1-2q_a)>0.
\]
All factors in these comparisons are strictly positive, by
Lemma~\ref{alg:bloch}. As $UKU^*=2T/\operatorname{tr}T$,
\[
 (Z^*KZ)_{ij}=\frac{T_{ab}}{\operatorname{tr}T\sqrt{p_ap_b}}
       \quad(i\in a,\ j\in b).
\]
This proves the phase assertion even for distinct indices in the same
block. The covariance and eigenvector identities follow from
$J^*J=I$ and $BU=0$. Thus
\[
 \Gamma_K(Z)=\operatorname{tr}(S Z^*KZ)
 =\operatorname{tr}(KZSZ^*)=N^2\lambda.
\]
In a neighborhood of $Z$ all weighted overlaps remain nonzero.
If $C_K(Y)$ is their phase matrix, differentiation with the real
Frobenius inner product gives $\nabla\Gamma_K(Y)=2KYC_K(Y)$.
At $Z$ this is $2KZS=2N\lambda KZ$, exactly the gradient of the
constraint term in \eqref{alg:Lagrangian}.
\end{proof}

\begin{proposition}[Exact balancing]\label{alg:balancing}
Among admissible integer four-weight profiles with total $N$,
$\lambda_p$ is largest and $\mu_p$ is smallest precisely when the four
counts differ by at most one. For every $N\geq10$ the balanced
profile is admissible. If $N\equiv2\pmod4$, then
\begin{equation}\label{alg:residue-two}
 \lambda_{\rm bal}=\frac{1+\sqrt{3-8/N^2}}4,\qquad
 \mu_{\rm bal}=\frac4{N+\sqrt{3N^2-8}-4}.
\end{equation}
\end{proposition}
\begin{proof}
For fixed $\lambda>2/3$, the function
$x\mapsto(\lambda-2x)^{-1}$ is strictly convex. If two counts differ
by at least two, transfer one from the larger to the smaller.
All counts stay positive and below $N/3$, and the sum of reciprocal
denominators strictly decreases. The coefficient $\lambda(3\lambda-2)$
in \eqref{alg:Phi} is positive, so at the old root the new $\Phi$ is
negative. Its unique root is strictly larger. Iteration terminates
at balanced counts, proving also the equality characterization.
For $N\geq10$, $\lceil N/4\rceil<N/3$.
For residue two the weights are $(a,a,b,b)$ with
$a+b=1/2$, $ab=1/16-1/(4N^2)$, and
\[
 \sum_a\frac1{\lambda-2p_a}
       =\frac{4\lambda-2}{\lambda^2-\lambda+4ab}.
\]
The root equation reduces to $2\lambda^2-\lambda-4ab=0$;
its relevant root is $(1+\sqrt{1+32ab})/4$.
\end{proof}


\subsection{The candidate domain needed at the endpoint}\label{app:candidate:domain}
Write $p_a=1/4+\varepsilon_a$, $e=\max|\varepsilon_a|$ and
$\sigma^2=\sum\varepsilon_a^2$. Throughout this subsection $e\le1/120$.
This includes every count profile in Proposition~\ref{prop:localization}
for $N\ge180$. All estimates concern the same alignment of the candidate
and a reference tetrahedron.

The exact scalar identities are
\[
 \Phi_p(\alpha)=4\sum_a\frac{\varepsilon_a^2}{\alpha-1/2-2\varepsilon_a},
 \qquad \Phi'_p(t)=4\left[2+\sum_a\frac{p_a(1-3p_a)}{(t-2p_a)^2}\right].
\]
For $f_t(p)=p(1-3p)/(t-2p)^2$,
\[
 f_t''(p)=\frac{2[t(4-3t)+4p(1-3t)]}{(t-2p)^4}>0
\]
on $2/3\le t\le\alpha<41/60$, $29/120\le p\le31/120$:
the numerator bracket is at least $43/200$. Jensen gives
$\Phi'_p(t)>4568/121>37$, whereas
$4/(\alpha-1/2-2e)<1200/49<74/3$. Therefore
\begin{equation}\label{app:candidate:deficit}
 0\le\alpha-\lambda_p\le\frac23\sigma^2,\qquad
 \frac{109}{100}<b:=3(2\lambda_p-1)<\frac{11}{10},\qquad
 u_0:=\frac34-\lambda_p>\frac1{15}.
\end{equation}
For balanced counts, $N^2\sigma^2=r(4-r)/4$. The odd budget is thus
$1/2$. For residue two the exact formula gives
$N^2(\alpha-\lambda_{\rm bal})=2/(\sqrt3+\sqrt{3-8/N^2})<11/19$
at $N\ge180$. These are the budgets in \eqref{eq:budget}.

\subsection{Exact projector and simultaneous polar estimates}
Let $\Pi=UU^*$ project onto $\ker B$. The two positive eigenvalues of $B$
have sum $4\lambda-1$ and product $D=3\lambda(2\lambda-1)$, whence
\[
 \Pi=I-\frac{(4\lambda-1)B-B^2}{D},\qquad
 \Pi_{aa}=\frac{3\lambda-2+2p_a}{b}.
\]
For complementary pairs $\{a,b\},\{c,d\}$, completion of $U$ to a
unitary matrix and equality of complementary minors give
\[
 \det\Pi[\{a,b\}]=\det(I-\Pi)[\{c,d\}]
 =\frac{\det B[\{c,d\}]}{D}=\frac{\lambda-1+p_a+p_b}{b}.
\]
These identities determine both the lengths $r_a=\|z_a\|$ and the
Bloch Gram matrix $Y_p$ of the normalized candidate spinors, relative
to the tetrahedral Bloch Gram matrix
$Y_0=\frac43I_4-\frac13\mathbf1\mathbf1^T$:
\begin{align}
 r_a^2&=\frac{1+4\varepsilon_a/b}{1+4\varepsilon_a},\label{app:candidate:projector}\\
 (Y_p-Y_0)_{ab}
 &=\frac{-\frac8{3b}(\varepsilon_a+\varepsilon_b)
                +\frac{64}{3b^2}\varepsilon_a\varepsilon_b}
 {(1+4\varepsilon_a/b)(1+4\varepsilon_b/b)}\quad(a\ne b).\label{app:candidate:Ygram}
\end{align}
The diagonal of $Y_p-Y_0$ is zero. The zero-sum identities imply
\[
 G:=\|Y_p-Y_0\|_F\le
 \frac{16/(3b)+112e/(3b^2)}{(1-4e/b)^2}\,\sigma
 \le\frac{11}{2}\sigma\le\frac{11}{120}.
\]
The coefficient at $b=109/100,e=1/120$ is $551200/100489<11/2$.

Here is a polar comparison used also for the actual group means.
For a real $3\times4$ unit-column matrix $X$ and a tetrahedron $V$ with
$V\mathbf1=0$, $VV^T=4I/3$, set $C_0=3XV^T/4$, $c=X\mathbf1/4$.
Then $X=C_0V+c\mathbf1^T$ and, if $G=\|X^TX-V^TV\|_F$,
\begin{equation}\label{app:candidate:polaridentity}
 G^2=\frac{16}{9}\|C_0^TC_0-I\|_F^2
       +\frac{32}{3}\|C_0^Tc\|^2+16\|c\|^4.
\end{equation}
This follows by expanding and using $V\mathbf1=0$. Also
$s_{\min}(C_0)^2\ge1-3G/4$. For the orthogonal polar factor $R$,
\begin{equation}\label{app:candidate:polar}
 \|X-RV\|_F^2\le P_G G^2,\qquad
 P_G=\max\left\{\frac{3/4}{(1+\sqrt{1-3G/4})^2},
                      \frac{3/8}{1-3G/4}\right\}.
\end{equation}
Indeed compare the two nonnegative terms
$\frac43\|C_0-R\|_F^2+4\|c\|^2$ with the first two terms of
\eqref{app:candidate:polaridentity}. At $G\le11/120$,
$s_{\min}(C_0)^2\ge149/160$, so $P_G<41/100$.
An orientation reversal is handled by conjugating the reference spinors.

For unit spinors $v,f$ and their Bloch vectors $x,y$,
\begin{equation}\label{app:candidate:spinor}
 d_{\rm ph}(v,f)^2=2(1-|v^*f|)
 =\frac{\|x-y\|^2}{2(1+|v^*f|)}.
\end{equation}
The total bound $(41/100)(11/120)^2<4[1-(999/1000)^2]$ allows one
common polar alignment and phases with $|v_a^*f_a|>999/1000$.
Moreover
\begin{gather*}
 |r_a^2-1|\le\frac{4e}{11(1-4e)}\le\frac1{319},\qquad
 \frac{499}{500}<r_a<\frac{501}{500},\\
 |r_a-1|\le c_r|\varepsilon_a|,
 \qquad c_r=\frac{4/11}{(1-4e)(1+499/500)}.
\end{gather*}
All these assertions use the larger $1/120$ domain; no smaller-domain
polar floor is extrapolated.

For the actual four count profiles, the following explicit formulas
sharpen the simultaneous errors. Put $\varepsilon_a=\delta_a/N$,
$\sigma=\sqrt C/N$, $b_0=109/100$, $k_0=8/(3b_0)$ and
$q_0=1-4e/b_0$. Write $\Delta_{\mathrm B}=Y_p-Y_0=D_0+F$, where
$(D_0)_{ab}=-8(\varepsilon_a+\varepsilon_b)/(3b)$ off the diagonal.
Exact subtraction in \eqref{app:candidate:Ygram} gives
\[
 F_{ab}=\frac{32}{3b^2}
 \frac{\varepsilon_a^2+4\varepsilon_a\varepsilon_b+\varepsilon_b^2
 +(4/b)(\varepsilon_a+\varepsilon_b)\varepsilon_a\varepsilon_b}
 {(1+4\varepsilon_a/b)(1+4\varepsilon_b/b)}.
\]
Thus $|F_{ab}|\le U_{ab}$, where
\[
 U_{ab}=\frac{32}{3b_0^2N^2q_0^2}
 \left(|\delta_a^2+4\delta_a\delta_b+\delta_b^2|
       +\frac4{b_0N}|(\delta_a+\delta_b)\delta_a\delta_b|\right).
\]
Let $f=(\sum_{a\ne b}U_{ab}^2)^{1/2}/\sigma$.
To see explicitly how this supplies a common polar bound, put
$E=9V\Delta_{\mathrm B}V^T/16$, $z=3V\Delta_{\mathrm B}\mathbf1/16$ and $P_0=(I+E)^{1/2}$.
The aligned Bloch error is
$H=(P_0-I)V+P_0^{-1}z\mathbf1^T$. Its linear parts satisfy
\[
 E_0=\frac{9k}{8}V\diag(\varepsilon_a)V^T,\quad
 z_0=-\frac{3k}{8}V\varepsilon,\quad k=8/(3b),
\]
\[
 \|E_0\|_F=\frac{3\sqrt2 k}{4}\sigma,
 \quad |z_0|=\frac{\sqrt3 k}{4}\sigma.
\]
These follow from $V^TV=\frac43I-\frac13\mathbf1\mathbf1^T$ and
$\sum\varepsilon_a=0$. Set
\[
 K_E=\frac{3\sqrt2}{4}k_0+\frac34f,\quad
 K_z=\frac{\sqrt3}{4}(k_0+f),\quad p_0=\sqrt{1-K_E\sigma},
\]
\[
 R_P=\frac{K_E^2\sigma}{2(1+p_0)^2},\qquad
 R_I=\frac{K_EK_z\sigma}{p_0(1+p_0)}.
\]
The scalar square-root identity and $P_0^{-1}-I$ give these full
nonlinear remainders. The tetrahedral Gram identities, applied separately
to each column, then give
\begin{align*}
 |H_a|/\sigma&\le\frac{\sqrt{3-4\delta_a^2/C}}{b_0}
                          +\frac{13}{16}f+R_P+R_I,\\
 \|H\|_F/\sigma&\le\frac{\sqrt8}{b_0}
                       +\frac{\sqrt{15}}4f+\frac2{\sqrt3}R_P+2R_I.
\end{align*}
For clarity, the linear point norm equals
$\sqrt{3\sigma^2-4\varepsilon_a^2}/b$; its aggregate norm is
$\sqrt8\,\sigma/b$. The nonlinear $F$ contribution is bounded by
$13\|F\|_F/16$ per column and $\sqrt{15}\|F\|_F/4$ in aggregate,
using $E_F=9VFV^T/16$, $z_F=3VF\mathbf1/16$ and
$\|E_F\|_F\le3\|F\|_F/4$, $|z_F|\le\sqrt3\|F\|_F/4$.
The two arrays $E_FV/2$ and $z_F\mathbf1^T$ are orthogonal, giving
the aggregate squared coefficient $15/16$; the point coefficient is
$3/8+\sqrt3/4<13/16$. This explains the paired, rather than independently
maximized, column dependence.

With $\phi=(501/500)/[2(1+999/1000)]$ and
$R_*=(13/16)f+R_P+R_I$, define
\begin{align*}
 C_{\mathrm{pt},a}&=c_r^2\delta_a^2/C+
      \phi[\sqrt{3-4\delta_a^2/C}/b_0+R_*]^2,\\
 C_{\mathrm{pt}}&=\max_a C_{\mathrm{pt},a},\\
 B_p&=(1/4+e)\left\{c_r^2+
 \phi[\sqrt8/b_0+(\sqrt{15}/4)f+(2/\sqrt3)R_P+2R_I]^2\right\}.
\end{align*}
In this same alignment they prove
\begin{equation}\label{app:candidate:PB}
 \max_a\|z_a-f_a\|^2\le C_{\mathrm{pt}}C/N^2,
 \qquad \sum_a p_a\|z_a-f_a\|^2\le B_p C/N^2.
\end{equation}
All divisions are guarded by $p_0>0$. On the full domain the generic
bound $f<61/100$ implies $K_E<33/8$, $K_z<77/32$, hence positivity;
actual profile bounds are sharper. Both signs and every permutation
are included in the absolute polynomials $U_{ab}$.

Finally the stress norm has a basis-free estimate. Let $j_a=1/\ell_a$
and $H=U^*\diag(j_a)U$. The identity $TL^{-1}\Pi=\Pi$ gives
$U^*TU=H^{-1}$ and $\|K-I\|=\operatorname{gap}(H)/\tr H$.
Writing $j_c=j(1/4,\lambda)$, exact division gives
\[
 j_a/j_c=1+C_0\varepsilon_a+
             C_0\varepsilon_a^2/(u_0-\varepsilon_a),\quad
 C_0=\frac{3\lambda-2}{(\lambda-1/2)u_0},\quad0\le C_0\le4.
\]
For $E=U^*\diag(\varepsilon_a)U$, the projector minors give
\[
 \tr E=2\sigma^2/b,\qquad
 \operatorname{gap}(E)^2=\sigma^2/3+
             4\sum\varepsilon_a^3/b+4\sigma^4/b^2.
\]
The remaining compression is positive semidefinite and has spectral
gap at most $4e^2/(u_0-e)$, whereas $\tr(H/j_c)\ge2$. Consequently
\begin{equation}\label{app:candidate:kappa}
 \kappa:=\|K-I\|\le
 2\sigma\sqrt{1/3+4e/b_0+4\sigma^2/b_0^2}
                         +\frac{2e^2}{1/15-e}<\frac{11}4e.
\end{equation}
The last endpoint check is
$4\sqrt{39013/106929}+2/7<11/4$. The radial and spinor estimates also give a generic bound $\eta$
for the candidate point error $\max_a\|z_a-f_a\|$:
$\eta^2\le[500/(319\cdot999)]^2+
\phi(41/100)(11/120)^2<1/1024$.
The profile-dependent choice of $\eta$ is specified in
\eqref{app:local:polar-input}.
Thus the candidate point error and $\kappa$ are below $1/32$.
For a fixed count profile all
analytic upper coefficients above improve as $N$ increases. Freezing
valid endpoint enclosures therefore proves the whole subsequent range;
no claim about monotonicity of rounded program outputs is needed.

\section{Localization of every column at the endpoint 180}\label{app:global}
This section proves the global and fine inputs to
Proposition~\ref{prop:localization}. All partitions below contain every
original column. The same partition is kept during the forward estimates.

\subsection{Deficit identity and estimates before partitioning}
Let $t_{ij}=(1+x_i\cdot x_j)/2$. The polynomial
\[
 f(t)=\alpha+\tfrac34(t-\tfrac12)-k(t^2-t+\tfrac16)
\]
satisfies $f(t)-\sqrt t=g(t)$ exactly. Put
$T_1=\sum_{ij}r_ir_jt_{ij}$. Tightness gives
$\sum_i r_i^2t_{ij}=N/2$, so
\[
 N^2/2-T_1=\tfrac12\sum_{ij}(r_i-r_j)^2t_{ij}\ge0.
\]
Also
$\sum_{ij}r_ir_j(t_{ij}^2-t_{ij}+1/6)
=\frac14\|\sum r_ix_ix_i^T-(L/3)I\|_F^2$.
Summing the polynomial identity proves
\begin{equation}\label{app:global:deficit}
 \alpha N^2-\Gamma_I(Y)=u+v+w+\tfrac34(N^2/2-T_1).
\end{equation}
Every term on the right is nonnegative. In particular $L\le N$.

Set $c=13/4$. On $N\ge180$, $u+v+w\le M\le11/19$ gives
$L/N>9999/10000$, using $a>77/250$. Since
$(77/250)(19999/10000)>8/13$,
\begin{equation}\label{app:global:radial}
 D_r:=\sum(r_i-1)^2=\frac{2u}{a(N+L)}\le\frac{cu}{N},
 \quad s=\sqrt{cu/N}<\frac{103}{1000}.
\end{equation}
Thus every radius is positive. With $E=\sum_{ij}g(t_{ij})$,
\begin{equation}\label{app:global:energy}
 E\le\frac{w}{(1-s)^2}\le\frac{M}{1-cM/N}<\frac{293}{500}.
\end{equation}
The second bound retains $u+w\le M$: after setting $x=\sqrt{cu/N}$,
$m=cM/N$, its numerator difference is
$m(1-x)^2-(m-x^2)(1-m)=(x-m)^2$.

Put $R_x=\sum x_i$, $M_0=\sum x_ix_i^T$, and
$l_i=(1,\sqrt3x_i)^T/2\in\mathbb R^4$. Cauchy--Schwarz and
$\sum r_i^2x_i=0$ give
\[
 \|M_0-NI/3\|_F\le2\sqrt{v/k}+\sqrt{2cu/3},\quad
 \|M_0-NI/3\|_{\rm op}<23/10,
\]
\[
 \|R_x\|^2\le(N/3+23/10)(2+s)^2D_r.
\]
For the operator estimate, a traceless symmetric $3\times3$ matrix
has squared operator norm at most two thirds of its squared Frobenius
norm, and
$(2/3)(1000/87+13/6)M<(23/10)^2$. The exact lifted identity is
\[
 \left\|\sum l_il_i^T-NI_4/4\right\|_F^2
 =\tfrac38\|R_x\|^2+\tfrac9{16}\|M_0-NI_3/3\|_F^2
 \le T(u,v):=Au+Bv+J\sqrt{uv},
\]
where
\begin{equation}\label{app:global:moment}
 A=\frac{19739570097}{6400000000},\quad B=\frac{375}{58},
 \quad J=\frac{1123}{200},\quad T(u,v)\le\frac{81}{10}(u+v).
\end{equation}
The last quadratic comparison has positive diagonal entries and determinant
\[
 \frac{293253687011}{928000000000}>0.
\]
If $z=u+v$, the traceless
$4\times4$ spectral inequality therefore yields
\begin{equation}\label{app:global:spectral}
 \lambda_{\min}\Bigl(\sum l_il_i^T\Bigr)
 \ge\Lambda_N(z):=N/4-\sqrt{243z/40}>0.
\end{equation}
Indeed an eigenvalue $t$ of a traceless $4\times4$ real symmetric matrix
forces the other squared eigenvalues to sum to at least $t^2/3$.

\subsection{Simultaneous anchors and coverage}
Let $\mathcal L=[l_1\ \cdots\ l_N]$ and $F=\mathcal L\mathcal L^T$.
Give each four-element subset $A$ probability
$\det(\mathcal L_A)^2/\det F$. For a set $J$ of size one or two,
\begin{equation}\label{app:global:inclusion}
 \Pr(J\subset A)=\det(\mathcal L_J^TF^{-1}\mathcal L_J).
\end{equation}
To verify this identity, whiten $U=F^{-1/2}\mathcal L$, so $UU^T=I$.
Cauchy--Binet expands
$\det(U\diag(t_i)U^T)=\sum_{|A|=4}\det(U_A)^2\prod_{i\in A}t_i$.
Taking the derivatives indexed by $J$ at $t_i=1$, and alternatively
expanding $\det(I+U\diag(t_i-1)U^T)$, gives
\eqref{app:global:inclusion}.

Write $q(t)=(3t-1)/2$, so $l_i^Tl_j=q(t_{ij})$.
Loewner order in the one- and two-point determinants gives
\[
 \Pr(i\in A)\le\Lambda_N(z)^{-1},\qquad
 \Pr(i,j\in A)\le[1-q(t_{ij})^2]\Lambda_N(z)^{-2}.
\]
The scalar bounds
$q(t)^2\le24g(t)$ for $t<2/3$ and
$1-q(t)^2\le96g(t)$ for $t\ge2/3$ follow from the canceled ratios
below, their derivatives and their values at $2/3$.
Explicitly, $g(t)/q(t)^2$ is at least its value at $2/3$, which exceeds
$1/24$, on the lower interval. On the upper interval
$g(t)/(1-t)>1/32$ and
$1-q(t)^2=(3/4)(1-t)(1+3t)\le3(1-t)$. Endpoints with zero factors
are handled directly or by continuity.
If $E_A$ is the sum of selected angular row energies, $H$ counts
selected unordered pairs with $t\ge2/3$, and
$X=2\sum_{\{i,j\}\subset A,t_{ij}<2/3}q(t_{ij})^2$, then
\[
 \mathbb E E_A\le E/\Lambda_N(z),\quad
 \mathbb E X\le24E/\Lambda_N(z)^2,\quad
 \mathbb E H\le48E/\Lambda_N(z)^2.
\]
Let $K_A=441/100$, $\eta_A=8/25$. For $E>0$, the expected potential
$E_A/(K_AE/N)+X/\eta_A^2+H$ is at most
\begin{equation}\label{app:global:potential}
 \frac{N}{K_A\Lambda_N(z)}+
 \left(24/\eta_A^2+48\right)
 \frac{(M-z)/(1-\sqrt{cz/N})^2}{\Lambda_N(z)^2}<1.
\end{equation}
The strict inequality is certified on $256$ closed intervals covering
$0\le z\le M$ at $N=180$, as specified in
Section~\ref{app:certificates}. For each interval $[l,h]$, replace
$\Lambda$ by $45-\sqrt{243h/40}$ and the energy by
\[
 \min\{M/(1-cM/180),(M-l)/(1-\sqrt{ch/180})^2\}.
\]
Every denominator is positive. This is a closed interval proof, not a
sampled maximization. Thus one actual four-subset has $H=0$, $x=\sqrt X<\eta_A$
and $E_A<(K_AE/N)(1-x^2/\eta_A^2)$.

The selected lifted Gram has smallest eigenvalue at least
$1-\sqrt3x/2$. The exact identity
\[
 \frac2{1+\sqrt{1-h}}(1-y)-(1-y^2/h)
       =\frac{(y-1+\sqrt{1-h})^2}{h}\ge0
\]
gives the coverage and disjointness ratios
\[
 R_{\rm cov}=\frac2{1+\sqrt{1-3\eta_A^2/4}},\qquad
 R_{\rm dis}=\frac2{1+\sqrt{1-\eta_A^2/2}}.
\]
For $h=\sqrt t$ and $b=1/\sqrt3$,
\[
 \frac{g(t)}{q(t)^2}=\frac{4(1-h)(3/4+kh)}{9(h+b)^2},\qquad
 \frac{g(t)}{1-t}=\frac{(h-b)^2(3/4+kh)}{1+h}.
\]
The first ratio decreases, with derivative $-4\sqrt3/[9(h+b)^3]$;
the second increases for $h\ge b$. The exact endpoint checks are
\[
 R_{\rm cov}\frac{K_A(293/500)}{180}
   <\frac{g(173/200)}{q(173/200)^2},\qquad
 R_{\rm dis}\frac{K_A(293/500)}{180}<\frac1{50}.
\]
If a data column were outside every $173/200$ cap, its unit lifted vector
would give $\sum_a q(t_{iA_a})^2\ge1-\sqrt3x/2$; its incident energies
would then contradict the first check. Two $83/100$ memberships of an
actual data column would imply an anchor squared distance at most
$(400/3)E_A$, using $g(t)/(1-t)>3/50$ there, whereas the selected-pair
bound gives at least $(8/3)(1-x/\sqrt2)$. The second check excludes it.
These assertions concern the data columns; the continuous caps need
not be disjoint. Since $173/200>83/100$, they give one partition into
four nonempty groups.

Let $\bar x_a$ be the means and
$S=\sum_a\sum_{i\in a}\|x_i-\bar x_a\|^2$. A valid lower membership
ratio at $173/200$ is
\[
 m_0=\frac{667855014484621414278044848317828318861}
 {9650268809435000000000000000000000000000}.
\]
Since means minimize squared distance,
$S\le4E_A/m_0<4K_AE/(m_0N)<255E/N$. We may weaken this to
$S\le264E/N$ for the subsequent common bootstrap.
If $E=0$, all $t_{ij}$ belong to $\{1/3,1\}$.
The positive lifted frame bound supplies four spanning data columns;
their lifted vectors are orthonormal. Every other unit lifted vector
has inner products in $\{0,1\}$ with that basis, and is one of them.
This gives the same partition and $S=0$ without division by $E$.

\subsection{Counts before fine localization}
Let $\Sigma_a$ be the within-group covariance, $S_a=\tr\Sigma_a$,
$\delta_a=m_a-N/4$, and let the square matrix $A_0$ have columns
$\sqrt{m_a}(1,\sqrt3\bar x_a)^T/2$. If
$F_0=\sum l_il_i^T-NI/4$, $H_0=\diag(0,\sum\Sigma_a)$ and
$t_0=(N-3S/4)/4$, then
\[
 A_0A_0^T-t_0I=F_0-\tfrac34(H_0-SI/4),\quad
 \diag(A_0^TA_0-t_0I)=\delta-\tfrac34(S_a-S/4)_a.
\]
The two Gram matrices have identical spectra. Their Frobenius norms,
diagonal contraction, and the centered PSD estimates imply
\begin{equation}\label{app:global:count}
 \sqrt C\le\sqrt{T(u,v)}+\frac{3\sqrt3}{4}S
                     \le\sqrt{T(u,v)}+\frac{13}{10}S.
\end{equation}
Each centered norm is at most $\sqrt3S/2$; both variance terms are
necessary. With $S\le264E/N$ and the shared budget, $256$ closed
$z$-interval checks give $C<5$ for $M=11/19$ and $C<19/4$ for $M=1/2$.
This precedes, and therefore does not rely on, any fine estimate.

For residue two the $\delta_a$ are half-integers with zero sum. The
strict square bound allows only $|\delta_a|\le3/2$, with at most one
entry of that magnitude, giving the two listed profiles. For odd $N$,
after sign reversal let $4\delta_a\equiv-1\pmod4$. Their squared sum
is below $76$. Magnitudes at least $9$ are excluded; an entry $7$
requires other entries summing to $-7$, whose minimum squares are
$25+1+1$, giving total $76$. Thus only $-5,-1,3$ remain. Their zero
sum permits either no $-5$, giving $(-1,-1,-1,3)$, or one $-5$,
giving $(-5,-1,3,3)$. Two $-5$ require two entries $3$ summing only
to $6$, so cannot have zero sum. This proves the complete count list.

\subsection{Shared-budget forward refinement}
Assume $S\le c_sE/N$ has already been proved on this partition. For a
closed region $u\le U,w\le W,u+v+w\le M$, let $\bar E$ be a valid
energy cap and set
\[
 p=1/4-m_c/N,\quad p_+=1/4+m_c/N,\quad
 S_*=c_s\bar E/N,\quad q=\sqrt{1-S_*/(pN)},\quad
 T_*=2S_* /(1+q),
\]
where $m_c=\max|\delta_a|$. All square roots and divisions require
their displayed positive domains. A sharp cap $\bar E$ follows by
maximizing $\min(W,M-u)/(1-\sqrt{cu/N})^2$ over $0\le u\le U$.
On the branch $w=M-u$ the sole interior maximum is
$u=(c/N)M^2$, of value $M/(1-cM/N)$; otherwise use the boundary.

Write
$Q=\sum_{a\ne b}m_am_b(\bar x_a\cdot\bar x_b+1/3)^2$.
The off-diagonal part of the preceding square mean-Gram calculation
gives $\sqrt Q\le(4/3)\sqrt T+S$. Write
$c_\times=\max_{a(i)\ne a(j)}(x_i\cdot x_j+1/3)$. Normalizing the
means gives the following three bounds, kept under a single budget supremum:
\begin{align}\label{app:global:finecaps}
 c_\times&\le\frac{4\sqrt T}{3Np}+b_s\sqrt E+d_sE,\\
 G&\le\frac{4\sqrt T}{3Np}+d_sE,\qquad
 \frac{\sqrt Q}{pN}\le\frac{4\sqrt T}{3Np}+d_mE,\nonumber\\
 b_s&=\sqrt{\frac{4c_s}{N(1+q)}},\quad
 d_m=\frac{c_s}{N^2p},\quad d_s=d_m\left(1+\frac{\sqrt6}{1+q}\right).
 \nonumber
\end{align}
Here $G$ bounds the Frobenius Gram error of the four normalized means.
The normalization costs follow from
$\sum_i\|x_i-\bar x_a/\|\bar x_a\|\|^2\le2S/(1+q)$.

On a closed bin $z\in[z_l,z_h]$, use
$u\le\min(U,z_h)$ and $w\le\min(W,M-z_l)$ to bound the energy and
the \emph{whole} expressions \eqref{app:global:finecaps}.
The raw moment $Au+Bv+J\sqrt{uv}$ can be maximized on $u+v=z_h$
by its concave one-variable expression. Its unconstrained upper bound
is $z_h(A+B+\sqrt{(A-B)^2+J^2})/2$; if $U_0=\min(U,z_h)$ obeys
\[
 2U_0\le z_h,\quad J^2(z_h-2U_0)^2
       \ge4(B-A)^2U_0(z_h-U_0),
\]
the constrained maximum is the value at $U_0$. This exact squared
derivative test avoids an uncertain radical branch.

There is also a spherical improvement to the cross bound. If the unit
mean dot product is at most $\gamma=-1/3+G$ and the two squared chord
errors sum to $T_s<2$, concavity of $\arccos(1-t/2)$ gives an angle
error at most $4\arcsin\sqrt{T_s/8}$. Put
\[
 C_s=1-T_s+T_s^2/8,\quad
 S_s=\sqrt{2T_s(1-T_s/8)}(1-T_s/4).
\]
When $\gamma<C_s$, the pair dot product is at most
$\gamma C_s+\sqrt{1-\gamma^2}S_s$; otherwise use $1$.
Here $T_s=2c_sE/[N(1+q)]$. This retains all within-group residuals.
On each closed budget bin, enclose both the linear bound for $c_\times$
in \eqref{app:global:finecaps} and the spherical bound just obtained,
adding $1/3$ to the latter dot-product bound. If the angular conditions
fail, use the trivial spherical bound $4/3$. Take the smaller of the
two upper enclosures on each bin, with outward rational rounding,
then the maximum over all bins covering the region. This gives a
regional upper bound. It may be reduced by taking its minimum with
other independently valid upper bounds on the same region. Denote the
resulting bound by $\bar c_\times$, so $c_\times\le\bar c_\times$
throughout the region.

Set $h_W=\sqrt{1-S_*/2}$ and
$h_C=\sqrt{(2/3+\bar c_\times)/2}$. The canceled scalar kernels are
\[
 K_W=\frac{(h_W-1/\sqrt3)^2(3/4+kh_W)}{4(1+h_W)},\quad
 K_C=\frac{(1-h_C)(3/4+kh_C)}{4(h_C+1/\sqrt3)^2}.
\]
Use lower enclosures of these kernels, $h_W>1/\sqrt3$, $h_C<1$, and
\begin{equation}\label{app:global:covariance}
 \mu_c=2\left[\frac13-\frac{23/10+S_*}{N}
                   -\frac{p_+T_*}{2(1+q)}\right],\quad
 \beta=2pK_W+\mu_cK_C>0.
\end{equation}
Then
\begin{equation}\label{app:global:jointcov}
 E\ge K_CQ+\beta NS.
\end{equation}
For detail, the cross-pair squared-dot sum is exactly
\[
 Q+2\sum_a\tr\left[\left(\sum_{b\ne a}m_b\bar x_b\bar x_b^T\right)
 \Sigma_a\right]+\sum_{a\ne b}\tr(\Sigma_a\Sigma_b).
\]
The last term is nonnegative. The full covariance is at least
$(N/3-23/10-S)I$. The variance of
$x_i\cdot(\bar x_a/\|\bar x_a\|)\in[1-T_*/2,1]$ gives its removed
same-group term at most $T_*S_a/[2(1+q)]$.
Indeed if its mean is $r_a=\|\bar x_a\|$, the scalar variance is at
most $(T_*/2)(1-r_a)$ and
$m_a(1-r_a)=S_a/(1+r_a)\le S_a/(1+q)$.
These facts yield $\mu_cNS$ for the cross part. The within-group
identity $\sum_{i,j\in a}\|x_i-x_j\|^2=2m_aS_a$ supplies
$2pK_WNS$, proving \eqref{app:global:jointcov}.

Starting from $264$, the recorded closed-domain inequalities prove
\begin{equation}\label{app:global:chain}
 264\longrightarrow240\longrightarrow183\longrightarrow107
 \longrightarrow58\longrightarrow40\longrightarrow34
 \longrightarrow32\longrightarrow31.
\end{equation}
Every arrow checks $\beta>1/c_{\rm next}$, hence is a forward implication,
not an assumed fixed point. The uniform chain uses $m_c=3/2$ and $128$
closed budget bins per stage. Subsequent regional refinement starts
from $31$, uses the already established profile's $m_c$ consistently,
and uses $64$ closed bins. No new partition is selected.

\subsection{Fine polar and spinor interface}
At the final regional coefficient let
$h_m=\sqrt{(2/3+\sqrt Q/(pN))/2}$, using its certified upper cap,
and put $K_m=K_C(h_m)$. If
$\gamma\le1/(4h_C^3)-2k$ is positive, then
$g(t)-\gamma t^2/2$ is convex on $[0,h_C^2]$. Indeed
$g^{\prime\prime}(t)=1/(4t^{3/2})-2k\ge\gamma$ for $t>0$,
and continuity extends convexity to $0$. Equivalently, $g$ is
$\gamma$-strongly convex on this interval. Jensen's inequality gives
\[
 E\ge K_mQ+[2pK_W+\mu_c\gamma/8]NS.
\]
Take a convex combination with \eqref{app:global:jointcov}, obtaining
$E\ge c_QQ+c_SNS$. If the positivity guard fails, use the first
inequality alone. The factor $1/8$ is the product of the strong-convexity
factor $1/2$ and the change of variable $t=(1+x\cdot y)/2$ squared.

Use the single polar alignment of the normalized means, with $P_G$ from
\eqref{app:candidate:polar} and $G<1/12$. The mean/residual decomposition
and Young's inequality with $\tau=201/200$ give
\[
 \sum_i\|x_i-v_{a(i)}\|^2\le A_SS+A_QQ/N,
\]
\[
 A_Q=p_+P_G\tau/p^2,\qquad
 A_S=\frac2{1+q}+
 \frac{6p_+P_G(1+1/(\tau-1))S_*}{p^2N(1+q)^2}.
\]
Therefore $L_B=\max(A_Q/c_Q,A_S/c_S)$ proves a total Bloch error
at most $L_BE/N$. This uses one joint covariance inequality, rather
than independently maximizing $Q$ and $S$.

Choose a phase floor $f$ from $127/128,63/64,31/32$ only if all three
strict guards hold:
\begin{equation}\label{app:global:phaseguards}
 L_B\bar E/N<4(1-f^2),\quad T_*<4(1-f^2),\quad
 P_GG^2<4(1-f^2).
\end{equation}
They concern, respectively, actual-to-polar, actual-to-normalized-mean
and normalized-mean-to-polar spinors. With $D_f=1/[2(1+f)]$, the
exact identity \eqref{app:candidate:spinor} and the quotient triangle
inequality give
\begin{align}\label{app:global:LU-LP}
 L_U&=D_fL_B,\\
 L_P&=\left[
 \sqrt{\frac{2D_f}{(1+q)\beta}+\frac{D_fP_G}{p^2NK_C}}
 +\frac{\sqrt{D_fP_G}\sqrt6\sqrt{\bar E/N^3}}
 {p(1+q)\beta}\right]^2.\nonumber
\end{align}
For the point estimate, the actual-to-mean contribution is bounded by
$2S/(1+q)$; the mean-to-polar contribution uses
$G\le\sqrt Q/(pN)+\sqrt6S/[pN(1+q)]$.
Maximizing their joint square-root terms under
$K_CQ+\beta NS\le E$ gives exactly the first square root in
\eqref{app:global:LU-LP}; the remaining normalization term gives the
second. Thus no separate maximum of $S$ and $Q$ is used in the first term.

In one common preliminary alignment and choice of column phases,
\begin{align}\label{app:global:physical-inputs}
 N\|Y-F\|_F^2&\le A_R:=cU+L_UW\frac{1+s_U}{(1-s_U)^2},\\
 N\max_i\|y_i-f_{a(i)}\|^2&\le P_R:=cU+L_PW\frac{1+s_U}{(1-s_U)^2},
 \qquad s_U=\sqrt{cU/N}.\nonumber
\end{align}
Indeed $d_{\rm ph}(r_iu_i,f)^2=(r_i-1)^2+r_id_{\rm ph}(u_i,f)^2$.
The sharper angular inputs are
$\sum d_{\rm ph}(u_i,f)^2\le L_U\bar E/N$ and
$\max d_{\rm ph}(u_i,f)^2\le L_P\bar E/N$.
The actual closest-gauge transport, required for the local Hessian,
is proved in Section~\ref{app:local}; a convenient preliminary
alignment is not substituted for that gauge.

\subsection{Continuous range and zero budgets}
Every accepted region is closed and all its constants cover its whole
interior and boundary. For larger $N$, at fixed deficits the energy
cap decreases, $N/\Lambda_N(z)$ decreases, and $\Lambda_N(z)$ increases.
All anchor and count inequalities therefore propagate. In the fine
chain, freeze the endpoint count floors, coefficient choices, closed
bins and phase floors. The energy and scatter decrease, mean lengths
improve, and every positive $N^{-1}$ or $N^{-1/2}$ error decreases.
The spherical cap is monotone geometrically because decreasing the
allowed mean dot product or chord error shrinks the feasible angle
range. Its algebraic maximizing branch need not remain fixed.
Propagate the finite chain in order, then its regional refinement and
polar interface. If $E=0$, \eqref{app:global:jointcov} gives $S=Q=0$;
all zero-error cases are direct. There is no division by $u,v,w$ or $E$
in the fine or local interface. These facts justify freezing the saved
180 localization coefficients for the later 182 local comparison.

\section{Local comparison}
\label{app:local}

This section proves the local implication used after
Sections~\ref{app:candidate} and \ref{app:global}.
All four profiles use the intrinsic residual and physical normal estimates.
For the odd profiles the complete unsigned comparison suffices.
For $C=3$ we use the stress-coordinate mean estimate with the complete
unsigned remainder. For $C=1$ we use the refined row coefficients and
the minimum of the complete signed alternatives and the complete unsigned
bound.
All derivatives are real derivatives on complex vector spaces, and
$\langle H,J\rangle=\Rea\tr(H^*J)$. Matrix norms without subscripts are
operator norms. Frame distances are measured in the Frobenius metric on
the original variables.

\subsection{Inputs, gauge, and the normal decomposition}
\label{app:local:inputs}

Let $a(i)\in\{1,2,3,4\}$ be the partition supplied by global localization,
let $m_a=|\{i:a(i)=a\}|$, and put $p_a=m_a/N$.
The candidate $Z=(z_{a(i)})_i$ and its positive stress $K$ satisfy
\begin{equation}\label{app:local:stationary-data}
 ZZ^*=\frac N2 I_2,\quad K>0,\quad\tr K=2,\quad
 S_{ij}=\frac{z_i^*Kz_j}{|z_i^*Kz_j|},\quad ZS=N\lambda Z.
\end{equation}
Every stress overlap is nonzero. Put $\alpha=(1+\sqrt3)/4$.
In one common alignment with unit tetrahedral spinors $f_a$, the candidate
estimates take the form
\begin{gather}
 e=\max_a|p_a-\tfrac14|,\qquad
 \sigma^2=\sum_a(p_a-\tfrac14)^2=C/N^2,\qquad
 \|K-I\|\le\kappa,\qquad 0\le\alpha-\lambda\le\tau,
 \label{app:local:errors}\\
 \max_a\|z_a-f_a\|\le\eta,\qquad
 \sum_a p_a\|z_a-f_a\|^2\le B\sigma^2,\qquad
 \eta=\sqrt{C_{\mathrm{pt}}C}/N.                                  \label{app:local:polar-input}
\end{gather}
Here $B,C_{\mathrm{pt}}$ are the simultaneous polar bounds of
Section~\ref{app:candidate}; the separate point and aggregate estimates
must use the same alignment. The four surviving profiles are
\[
\begin{array}{c|cccc}
 C&3/4&11/4&1&3\\ \hline
 m_*=\max_a|m_a-N/4|&3/4&5/4&1/2&3/2\\
 L_*=\sum_a|m_a-N/4|&3/2&3&2&3
\end{array}
\]
with both signed profiles and all permutations included.
Throughout the certified range, $e\le1/120$ and $\eta,\kappa<1/32$.
Write $p_-=1/4-e$, $p_+=1/4+e$, and $\omega_p=1-4e>0$.

The competitor is any normalized left kernel $Y$ with
$YY^*=NI_2/2$ and $YA=0$. There is no assumption that its columns have
equal lengths or take only four values. On a closed budget region,
global localization supplies nonnegative numbers $A_R,P_R,L_U,L_P,\bar E$
such that, in its preliminary alignment,
\begin{gather}
 \|Y-F\|_F^2\le A_R/N,\qquad
 \max_i d_{\rm ph}(y_i,f_{a(i)})^2\le P_R/N,
 \label{app:local:physical-input}\\
 \sum_i d_{\rm ph}(u_i,f_{a(i)})^2\le L_U\bar E/N,\qquad
 \max_i d_{\rm ph}(u_i,f_{a(i)})^2\le L_P\bar E/N,       \label{app:local:unit-input}
\end{gather}
where $u_i=y_i/\|y_i\|$ and $d_{\rm ph}$ minimizes over a scalar phase.
The first norm is evaluated in the preliminary phase choices.
The physical lengths lie in $[1-s,1+s]$, with $s<1$.
These bounds hold throughout the region.

For fixed $K$ define
\[
 \Gamma_K(Y)=\sum_{i,j}|y_i^*Ky_j|,\qquad
 \mathcal L(Y)=\Gamma_K(Y)-N\lambda\tr(KYY^*),\qquad
 Q(H,J)=\tfrac12D^2\mathcal L(Z)[H,J].
\]
When one argument is displayed, $Q(H)$ means $Q(H,H)$.

\begin{lemma}[Closest gauge]\label{app:local:gauge}
Choose a minimizer of $\|UYD-Z\|_F$ over $U\in U(2)$ and diagonal
unitary $D$, and relabel that representative as $Y$. Set
$\Delta=Y-Z$, $d=\|\Delta\|_F$, and
\[
 C_n=\Delta\Delta^*/N,\qquad R_n=C_nZ,\qquad \Xi=\Delta+R_n .
\]
Then
\begin{gather}
 D\mathcal L(Z)=0,\qquad \Gamma_K(Z)=N^2\lambda,\qquad
 \Delta Z^*=Z\Delta^*=-\tfrac12\Delta\Delta^*,            \label{app:local:gauge-id}\\
 \Xi Z^*=0,\quad \Ima(z_i^*\Xi_i)=0,\quad
 \langle\Xi,R_n\rangle=0,\nonumber\\
 d^2=\|\Xi\|_F^2+\|R_n\|_F^2,\quad
 \|R_n\|_F/d\le d/\sqrt{2N}\quad(d>0).                 \label{app:local:normal-id}
\end{gather}
\end{lemma}
\begin{proof}
The orbit is compact. Differentiating its squared distance at a minimum
gives $\langle\Delta,AY\rangle=0$ for every skew-Hermitian $A$ and
$\langle\Delta,iY\Phi\rangle=0$ for every real diagonal $\Phi$.
The terms with $Y$ replaced by $\Delta$ vanish, so $\Delta Z^*$ is
Hermitian and $\Ima(z_i^*\Delta_i)=0$. Equal covariances give
$\Delta Z^*+Z\Delta^*=-\Delta\Delta^*$.
Multiplying the definition of $\Xi$ by $Z^*$ gives its horizontal
constraint. Since $C_n$ is Hermitian, its individual phase equations
also hold. Horizontal arrays are orthogonal to $C_nZ$, and
\[
 \|R_n\|_F^2=\frac{\tr((\Delta\Delta^*)^2)}{2N}
 \le\frac{d^4}{2N}.
\]
Finally $\nabla\Gamma_K(Z)=2KZS=2N\lambda KZ$ proves stationarity, and
$\Gamma_K(Z)=\tr(KZSZ^*)=N^2\lambda$.
\end{proof}

The simultaneous change of matrix is $A\mapsto D^*AD$:
$(UYD)(D^*AD)=UYA D=0$. Thus diagonal entries, rank and off-diagonal
moduli are preserved. We choose the representative first and then fix
its stress $K$; $\Gamma_K$ need not be invariant under common unitaries
when $K$ is fixed.
For equal covariances,
$\Gamma_K(Y)-\Gamma_K(Z)=\mathcal L(Y)-\mathcal L(Z)$.

Split $\Xi=H+U$ into cluster means and zero-sum within-cluster residuals.
This is an orthogonal splitting in the Frobenius metric.
Both $H$ and $U$ satisfy the horizontal and individual phase equations.
Every Hessian coefficient at $Z$ depends only on its two cluster types;
summing one residual factor gives zero. Consequently
\begin{equation}\label{app:local:split}
 Q(H,U)=Q(U,R_n)=0 .
\end{equation}

\subsection{The reference tensor and an endpoint perturbation bound}
\label{app:local:reference}

For a cluster-constant array, put
$\|H\|_\nu^2=\frac14\sum_a\|H_a\|^2$. More generally give each member
of cluster $a$ weight $\nu_i=1/(4m_a)$.
The normalized reference potential is
\[
 \ell_0(Y)=\sum_{i,j}\nu_i\nu_j|y_i^*y_j|
                  -\alpha\sum_i\nu_i\|y_i\|^2,\qquad
 B_0=D^2\ell_0(F),\quad q_0=B_0/2 .
\]
It has zero gradient at $F$: its phase matrix is
$\sqrt3 f_i^*f_j+(1-\sqrt3)\mathbf1_{a(i)=a(j)}$.
All infinitesimal common-unitary and individual-phase directions lie
in the full kernel of $B_0$, including their pairing with ambient vectors.

\begin{lemma}[Complete reference mean tensor]\label{app:local:reference-spectrum}
On
\[
 \mathcal H_0=\{H:\tfrac14\sum_aH_af_a^*=0,\
                         \Ima(f_a^*H_a)=0\},
\]
the full mean Hessian has eigenvalues $-2/\sqrt3$ with multiplicity three
and $-\sqrt3$ with multiplicity two. On traceless Hermitian normal
arrays $MF$, it equals $(2/\sqrt3)I$; the scalar normal is null.
The horizontal-to-normal operator norm is $\sqrt{2/3}$, and the full
ambient mean operator norm is at most $7/4$.
\end{lemma}
\begin{proof}
Let $v_a$ be tetrahedral Bloch vectors:
$\sum_av_a=0$, $\sum_av_av_a^T=4I_3/3$.
Every phase-horizontal mean has the unique representation
\[
 H_a=a_af_a+\tfrac12(u_a\cdot\sigma)f_a,\qquad
 a_a\in\mathbb R,\quad u_a\perp v_a,\qquad
 \|H_a\|^2=a_a^2+\|u_a\|^2/4.
\]
Here $\sigma$ denotes the Pauli matrices. Put
$A_1=\sum_aa_av_a$, $U_1=\sum_au_a$ and
$D_{ab}=u_a\cdot v_b+v_a\cdot u_b$.
Differentiating
$r_ar_b\sqrt{(1+x_a\cdot x_b)/2}$ at
$r_a=1$, $x_a=v_a$ gives
\begin{align}
 B_0(H,H)={}&-\frac1{2\sqrt3}\sum_aa_a^2
 +\frac{(\sum_aa_a)^2}{8\sqrt3}
 +\frac{\sqrt3}{16}A_1\cdot U_1
 +\frac{\sqrt3}{32}\|U_1\|^2
 -\frac{3\sqrt3}{256}\sum_{a\ne b}D_{ab}^2 .
 \label{app:local:reference-form}
\end{align}
For verification, along a spherical curve with initial velocity $u_a$,
the cross-pair square-root factor has derivatives
$D_{ab}/(4h)$ and
\[
 \frac{u_a\cdot u_b+(\|u_a\|^2+\|u_b\|^2)/6}{2h}
             -\frac{D_{ab}^2}{16h^3},\qquad h=1/\sqrt3.
\]
Summing ordered pairs with weight $1/16$ gives
\eqref{app:local:reference-form}; accelerations do not matter at a
stationary point.

The matrix constraint is equivalent to
$\sum a_a=0$, $U_1=-2A_1$, and $\sum u_a\times v_a=0$.
Set $b=3A_1/4$. The component
$a_a=b\cdot v_a$, $u_a=-(b-a_av_a)$ has squared $\nu$ norm
$\|b\|^2/2$ and Hessian $-\sqrt3\|b\|^2/3$.
Its orthogonal complement has $a_a=0$ and $U_1=0$.
Writing $V=(v_a)$ and $U=(u_a)$, the latter is $U=TV$ with $T$ symmetric,
and $D=2V^TTV$ gives $\|D\|_F^2=(16/3)\|U\|_F^2$.
In coordinates $v_a=(\pm1,\pm1,\pm1)/\sqrt3$ with an even number of
minus signs, $v_a^TTv_a=0$ makes $T$ diagonal and traceless.
This gives the two shape eigenvalues $-\sqrt3$.

For $H_n=(b\cdot\sigma)F$, substitution and polarization give
\[
 \|H_n\|_\nu^2=\|b\|^2,\quad
 B_0(H_n,H_n)=2\|b\|^2/\sqrt3,\quad
 B_0(H,H_n)=-\sqrt3\,b\cdot A_1/4 .
\]
The horizontal constraints imply
$\|H\|_\nu^2\ge9\|A_1\|^2/32$, with equality on the first sector.
Thus the mixed norm is $\sqrt{2/3}$. Its coupled blocks have eigenvalues
$\pm\sqrt2$, the shape eigenvalues have modulus $\sqrt3$, and all
remaining gauge/scalar modes are zero. Hence the ambient norm is
$\sqrt3<7/4$.
\end{proof}

\begin{lemma}[Squared reference projection]\label{app:local:projection}
Suppose a mean array $G$ satisfies
\[
 \left\|\tfrac14\sum_aG_af_a^*\right\|_F\le b\|G\|_\nu,
 \qquad |\Ima(f_a^*G_a)|\le\phi\|G_a\|.
\]
Its reference-horizontal and Hermitian-normal components satisfy
\[
 \|G_0\|_\nu^2\ge[1-\phi^2-3(b+\phi)^2]\|G\|_\nu^2,\qquad
 \|G_n\|_\nu\le\sqrt2 b\|G\|_\nu .
\]
Consequently, if $\beta=\phi^2+3(b+\phi)^2<1$,
\begin{equation}\label{app:local:reference-gap}
 B_0(G,G)\le-g(b,\phi)\|G\|_\nu^2,\quad
 g(b,\phi)=\frac2{\sqrt3}(1-\beta)
                         -\frac4{\sqrt3}(b+b^2).
\end{equation}
\end{lemma}
\begin{proof}
Remove the reference phase component, of norm at most $\phi\|G\|_\nu$.
The remaining constraint matrix has norm at most $(b+\phi)\|G\|_\nu$.
On phase-horizontal arrays the adjoint constraint is
$M\mapsto P_{\rm ph}(Mf_a)$, and for $M=A+iB$ with $A,B$ Hermitian,
$\tr B=0$,
\[
 \|P_{\rm ph}(MF)\|_\nu^2=\tfrac12\|A\|_F^2+\tfrac13\|B\|_F^2 .
\]
Thus its least nonzero singular value is $1/\sqrt3$.
The two successive corrections are orthogonal, proving the squared
horizontal estimate. If $C=\frac14\sum G_af_a^*$, the normal projection
is $2\operatorname{Herm}(C)F$ and has norm squared
$2\|\operatorname{Herm}(C)\|_F^2$. Reference gauge directions are
orthogonal to this normal and annihilated by $B_0$.
Lemma~\ref{app:local:reference-spectrum}, $\|G_0\|_\nu\le\|G\|_\nu$,
and the normal estimate yield \eqref{app:local:reference-gap}.
\end{proof}

The following endpoint version of the modulus derivative estimate is
useful because a certified overlap floor concerns actual overlaps,
not the intervening line segment.
\begin{lemma}[Endpoint Hessian bound]\label{app:local:endpoint-hessian}
For nonzero $u,v\in\mathbb C\simeq\mathbb R^2$, and
$0<h\le\min(|u|,|v|)$,
\[
 \|D|\cdot|(u)-D|\cdot|(v)\|\le |u-v|/h,\qquad
 \|D^2|\cdot|(u)-D^2|\cdot|(v)\|
                           \le\frac76\,|u-v|/h^2 .
\]
\end{lemma}
\begin{proof}
The case $u=v$ is immediate. Write $r=|u|\ge s=|v|$,
$t=r/s$, $d=t-1$, $c=\cos\arg(u\bar v)$, and
$D=|u-v|/s=\sqrt{d^2+2t(1-c)}$.
The gradient difference has squared norm $(D^2-d^2)/t\le D^2$.
The norm of the Hessian difference, multiplied by $s$, is exactly
\[
 \frac{d+\sqrt{d^2+4t(1-c^2)}}{2t}.
\]
Put $L=7/6$ and $y=d/D$. Since
$1-c^2\le(D^2-d^2)/t$,
\[
 tL^2D^2-dLD-(1-c^2)
 \ge\frac{D^2}{t}\left[(y-tL/2)^2+\frac34t^2L^2-1\right]>0.
\]
Also $2tLD\ge d$. Squaring in the correct direction proves the displayed
Hessian norm is at most $LD$. Rescaling and $s\ge h$ prove the claim.
No lower bound on intermediate segment moduli was assumed.
\end{proof}

For later use define, for a certified cross stress floor $0<h_c\le15/26$,
\begin{align}
 d_p&=2\eta+\eta^2+\kappa(1+\eta)^2,\quad
 u_p=\eta+\kappa(1+\eta), \notag\\
 T_p&=\frac{7d_p}{6h_c^2}
       +\frac{u_p(2+u_p)}{h_c}+\frac{d_p}{2h_c}+\frac\kappa2 .
 \label{app:local:pair-error}
\end{align}
Then the absolute \emph{full} cross-pair Hessian error is at most
$T_p(\|h_a\|+\|h_b\|)^2$. Indeed the overlap, first-variation and
second-variation errors are bounded respectively by $d_p$, $u_pv$,
and $\kappa v^2/2$, where $v=\|h_a\|+\|h_b\|$; the corresponding reference
variations have bounds $v$ and $v^2/2$.
Apply Lemma~\ref{app:local:endpoint-hessian} to the reference
first-variation arguments, then change those arguments, the gradient,
and the second variation. This gives precisely the four terms in
\eqref{app:local:pair-error}. Same-type \emph{mean} pairs have exact
Hessian $2h_a^*Kh_a$, and hence error at most $2\kappa\|h_a\|^2$.

\subsection{Mean curvature in physical, mass, and stress coordinates}
\label{app:local:mean}

An ambient physical mean perturbation bound needed for the normal terms is
\begin{align}
 W_\times&=\max_{a\ne b}|16p_ap_b-1|,\qquad W_d=8e+16e^2,\notag\\
 E_{\mathrm{Hess}}&=(1+4e)^2(3T_p+\kappa/2)
       +(27/4)W_\times+W_d/2 \notag\\
 &\quad+2\{(1+4e)[(7/10+\tau)\kappa+\tau]+(14/5)e\}.
 \label{app:local:ambient-P}
\end{align}
It satisfies
\begin{equation}\label{app:local:ambient-perturbation}
 |N^{-2}Q(H,J)-q_0(H,J)|
                  \le(E_{\mathrm{Hess}}/2)\|H\|_\nu\|J\|_\nu
\end{equation}
on all cluster-constant arrays, without a horizontal hypothesis.
To verify this, use
$\sum_{a\ne b}(l_a+l_b)^2/16\le3\|H\|_\nu^2$.
The reference cross Hessian is bounded by $(9/4)(l_a+l_b)^2$;
the same-type one is exactly $2l_a^2$. This gives the cross and diagonal
weight errors separately. The multiplier error follows from
$\|\lambda K-\alpha I\|\le(7/10+\tau)\kappa+\tau$ and
$|4p_a-1|\le4e$. Polarization of the resulting absolute self-adjoint
operator bound gives \eqref{app:local:ambient-perturbation}.
For the four profiles, respectively, valid cross-weight envelopes are
\[
 W_\times\le 2/N+3/N^2,\quad 6/N+15/N^2,\quad
                 4/N+4/N^2,\quad 4/N+12/N^2 .
\]
They follow directly from
$16p_ap_b-1=4(\delta_a+\delta_b)/N+16\delta_a\delta_b/N^2$.

\begin{lemma}[Mass coordinates]\label{app:local:mass-mean}
Let $\epsilon_+\ge\max_a(p_a-1/4)$ cover both signs of the count profile,
and define
\[
 b_g=\sqrt{B\sigma^2/\omega_p},\quad
 g_g=g(b_g,\eta)-(1+2/\omega_p)T_p-\kappa/2
                   -2(\alpha\kappa+\tau)/\omega_p-8\alpha\epsilon_+/\omega_p .
\]
If the reference projection fraction is positive and $g_g>0$, then
\begin{equation}\label{app:local:mass-gap}
 Q(H)\le-Nc_g\|H\|_F^2,\qquad c_g=\omega_p g_g/2
\end{equation}
for every actual horizontal mean.
\end{lemma}
\begin{proof}
Set $G_a=4p_aH_a$. Then $\frac14\sum G_az_a^*=0$ and
$\Ima(z_a^*G_a)=0$. The reference constraint error is at most
$b_g\|G\|_\nu$, and its phase error is at most $\eta$.
Lemma~\ref{app:local:projection} gives $g(b_g,\eta)$.
The exact norm relation is
$\|H\|_F^2/N=\frac14\sum\|G_a\|^2/(4p_a)\le\|G\|_\nu^2/\omega_p$.

For completeness, the mass cancellation can be checked at the
half-Hessian level. Put
$\theta^0_{ab}=f_a^*f_b/|f_a^*f_b|$ and
$A_{ab}(g)=[\Ima(\overline{\theta^0_{ab}}g^*f_b)]^2/|f_a^*f_b|$.
If $q_p^F$ denotes the reference half-Hessian with masses $p$, then
\[
 q_p^F(H)-q_{1/4}^F(G)
 =\sum_a\frac{\sum_b(p_b-p_a)A_{ab}(G_a)
                         +4\alpha(p_a-\tfrac14)\|G_a\|^2}{4(4p_a)} .
\]
Mixed first-variation and second-variation terms cancel because
$p_ap_b/[(4p_a)(4p_b)]=1/16$.
In coordinates $f_a=(1,0)$, $g=(x+iy,u+iv)$,
$\sum_bA_{ab}(g)=(1+\sqrt3)y^2+\sqrt3(u^2+v^2)
\le4\alpha\|g\|^2$.
The displayed difference is therefore at most
$4\alpha\epsilon_+\|G\|_\nu^2/\omega_p$.

Moreover
\[
 \sum_{a\ne b}p_ap_b(\|H_a\|+\|H_b\|)^2
                       \le(1+2/\omega_p)\|G\|_\nu^2 .
\]
Expanding the two squared-column terms gives at most
$(4/\omega_p-1)\sum\|G_a\|^2/8$, and the product terms at most
$3\sum\|G_a\|^2/8$.
The same-type full-Hessian parameter error is $\kappa\|G\|_\nu^2/2$,
and the covariance error is $2(\alpha\kappa+\tau)\|G\|_\nu^2/\omega_p$.
These estimates give $g_g$. Its strict positivity is checked before
converting the negative norm estimate.
\end{proof}

\begin{lemma}[Stress and mass coordinates]\label{app:local:stress-mean}
Put
\begin{gather*}
 d_K=\frac{\kappa}{1+\sqrt{1-\kappa}},\qquad
 \eta_v=\sqrt{1+\kappa}\eta+d_K,\qquad
 b_v=\frac{\sqrt{B\sigma^2}+d_K}{\sqrt{\omega_p}},\\
 \phi_v=\eta_v+\frac{\kappa(1+\eta)}{\sqrt{1-\kappa}},\qquad
 d_v=2\eta_v+\eta_v^2,\\
 T_v=\frac{7d_v}{6h_c^2}
               +\frac{\eta_v(2+\eta_v)}{h_c}+\frac{d_v}{2h_c},\\
 g_v=g(b_v,\phi_v)-(1+2/\omega_p)T_v-2\tau/\omega_p-8\alpha\epsilon_+/\omega_p .
\end{gather*}
Assume $\eta_v<1/32$, $\phi_v^2+3(b_v+\phi_v)^2<1$ and $g_v>0$.
Then
\begin{equation}\label{app:local:stress-gap}
 Q(H)\le-Nc_v\|H\|_F^2,\qquad c_v=\omega_p(1-\kappa)g_v/2 .
\end{equation}
\end{lemma}
\begin{proof}
Set $T=K^{1/2}$, $v_a=Tz_a$, and $G_a=4p_aTH_a$.
The entire Lagrangian in $TY$ has pair stress $I$ and multiplier
$\lambda I$. The transformed covariance is $K/2$, not $I/2$;
no tightness of $v$ is used.
Exactly $\frac14\sum G_av_a^*=0$, but its phase equation has the defect
\[
 |\Ima(v_a^*G_a)|
 =4p_a|\Ima(z_a^*(K-I)H_a)|
 \le\frac{\kappa(1+\eta)}{\sqrt{1-\kappa}}\|G_a\|.
\]
Spectral calculus gives $\|T-I\|\le d_K$, proving the point bound
$\|v_a-f_a\|\le\eta_v$. For the aggregate use
$v-f=(z-f)+(T-I)z$ and the \emph{original} identity
$\sum p_a\|z_a\|^2=1$:
\[
 \left(\sum p_a\|v_a-f_a\|^2\right)^{1/2}
                         \le\sqrt{B\sigma^2}+d_K .
\]
Thus Lemma~\ref{app:local:projection} applies with $b_v,\phi_v$;
in particular the transformed phase defect has not been erased.

The cross-pair perturbation is now
\eqref{app:local:pair-error} with point error $\eta_v$ and stress
error zero, namely $T_v$. The overlap floor is unchanged, since
$v_a^*v_b=z_a^*Kz_b$.
Same-type parameter error is exactly zero, because both kernels have
Hessian $2\|g_a\|^2$. Their mass error is still included in the preceding
mass-coordinate identity. The covariance error is $2\tau/\omega_p$.
This proves the Hessian bound with coefficient $g_v$.
Finally
\[
 \|G\|_\nu^2=\sum_a4p_a^2H_a^*KH_a
             \ge \omega_p(1-\kappa)\|H\|_F^2/N .
\]
Its substitution into a negative bound is valid after $g_v>0$.
The transformation is invertible and introduces no equality mode.
\end{proof}

The accepted mean coefficient may be any certified lower bound for
$\max(c_g,c_v)$, using only those alternatives whose guards hold.
The stress transform is confined to this mean estimate. In particular
neither $b_g$ nor $b_v$ replaces the physical normal projection below.

\subsection{The intrinsic residual bound}
\label{app:local:residual}

We use the two-dimensional spectral reduction, since a trace-only
bound would lose the anisotropy needed in the nonuniform profiles.
Use the candidate construction of Section~\ref{app:candidate}, writing
its auxiliary rank-two Gram matrix as $\mathcal B$ to distinguish it
from the aggregate coefficient $B$ in \eqref{app:local:polar-input}.
Define
\begin{gather}
 \ell_a=\frac{1-\lambda-p_a}{(1-\lambda)(\lambda-2p_a)},\quad
 L=\diag(\ell_a),\quad \mathcal T=L-L\mathcal B L,\quad
 \mathcal T_0=\tr\mathcal T, \notag\\
 d_a=\ell_a[1-\ell_a(\lambda-p_a)],\quad
 r_{Ka}^2=\frac{d_a}{p_a\mathcal T_0},\qquad
 s_{ab}=\frac{\ell_a\ell_b}{\mathcal T_0}\quad(a\ne b),\quad
 s_{aa}=r_{Ka}^2.                                    \label{app:local:intrinsic-scalars}
\end{gather}
All displayed denominators and $d_a$ are strictly positive.
If the columns of $U$ form an orthonormal basis of
$\operatorname{range}\mathcal T$, then
$K=2U^*\mathcal T U/\mathcal T_0$ and
$z_a=U^*e_a/\sqrt{2p_a}$. The exact identity
$\mathcal T L^{-1}\mathcal T=\mathcal T$ implies, for $v_a=K^{1/2}z_a$,
\begin{equation}\label{app:local:weighted-resolution}
 \sum_a(p_a/\ell_a)v_av_a^*=I_2/\mathcal T_0 .
\end{equation}
Indeed $U^*L^{-1}U=(U^*\mathcal T U)^{-1}$.

\begin{lemma}[Pointwise intrinsic form]\label{app:local:intrinsic}
For a physical vector $h$ at type $a$ define
\[
 q_a(h)=\sum_b\frac{p_b}{s_{ba}}
       [\Ima(\overline{\theta_{ba}}z_b^*Kh)]^2-\lambda h^*Kh,
 \qquad \theta_{ba}=z_b^*Kz_a/s_{ba}.
\]
For the full within-cluster zero-sum residual $U$,
$Q(U)/N=\sum_jq_{a(j)}(U_j)$.
For $j\ne a$ put
\[
 u_j=p_j^2(r_{Kj}^2-s_{aj}^2/r_{Ka}^2),\qquad
 D_a=\frac{2\sum_{j<k}u_ju_k-\sum_ju_j^2}{4},
\]
where the sums run over the three indices different from $a$, and set
\begin{gather}
 \mathcal D_a=D_a\sum_{\substack{j<k\\j,k\ne a}}
                       \frac1{p_jp_ks_{aj}s_{ak}},\qquad
 c_{Ka}=\lambda-\frac{\ell_a^{-1}
                   +\sqrt{\ell_a^{-2}-4\mathcal D_a}}2 .
 \label{app:local:intrinsic-eigen}
\end{gather}
If $c_K\le\min_a\{\lambda,c_{Ka}\}$ is strictly positive, then
\begin{equation}\label{app:local:physical-residual}
 q_a(h)\le-(1-\kappa)c_K\|h\|^2
 \quad\text{whenever }\Ima(z_a^*h)=0 .
\end{equation}
\end{lemma}
\begin{proof}
In the ordered Hessian, mixed-argument and mixed first-variation terms
of a zero-sum residual factor through its group sums and vanish.
The two pure-first terms give $Q(U)/N=\sum_jq_{a(j)}(U_j)$,
including same-cluster and diagonal pairs.

Fix $e_0=v_a/r_{Ka}$ and a perpendicular unit vector $e_1$.
Rotate $v_b$ into its pair phase, writing
$(s_{ab}/r_{Ka})e_0+c_be_1$, with $c_a=0$.
Stress stationarity gives
$\sum_bp_bc_b=0$ and $\sum_bp_bs_{ab}=\lambda r_{Ka}^2$.
If $K^{1/2}h=(x+iy)e_0+u e_1$, direct expansion therefore yields
\[
 q_a(h)=-\lambda x^2+
         \sum_{b\ne a}\frac{p_b}{s_{ab}}
                [\Ima(\bar u c_b)]^2-\lambda|u|^2 .
\]
All terms containing $y$ cancel exactly; the whole $iz_a$ direction
is in the kernel, including its bilinear pairings.

Let $M_a$ be the real two-dimensional positive quadratic form in the
middle sum. Each real rank-one form $[\Ima(\bar u c_b)]^2$
has trace $|c_b|^2$. Since $c_a=0$ and $s_{ab}=\ell_a\ell_b/\mathcal T_0$ for $b\ne a$,
\begin{align*}
 \tr M_a&=\sum_{b\ne a}\frac{p_b|c_b|^2}{s_{ab}}
 =\frac{\mathcal T_0}{\ell_a}\sum_b\frac{p_b}{\ell_b}|c_b|^2,\\
 \sum_b\frac{p_b}{\ell_b}|c_b|^2
 &=e_1^*\left(\sum_b\frac{p_b}{\ell_b}v_bv_b^*\right)e_1
 =\frac1{\mathcal T_0}.
\end{align*}
The last equality is \eqref{app:local:weighted-resolution}; pair-phase
rotations do not change the projected moduli. Thus $\tr M_a=1/\ell_a$.
The vectors $t_j=p_jc_j$ form a closed triangle. Their pairwise
determinants have common square $D_a$, since
$t_j\cdot t_k=(u_l-u_j-u_k)/2$.
The two-dimensional Cauchy--Binet formula gives $\det M_a=\mathcal D_a$.
This proves its \emph{complete} largest-eigenvalue formula
\eqref{app:local:intrinsic-eigen}.

For a physical phase-horizontal $h$, choose real $t$ so that
$h'=h-it z_a$ satisfies $\Ima(z_a^*Kh')=0$.
Then $q_a(h')=q_a(h)$, while
$\|h'\|^2=\|h\|^2+t^2\|z_a\|^2\ge\|h\|^2$.
The spectral bound and $K\ge(1-\kappa)I$ prove
\eqref{app:local:physical-residual}.
This is a pointwise statement, prior to imposing any zero group sum.
The phase adjustment is real-linear within each group, so it would also
preserve such sums if applied to an array.
\end{proof}

In the accepted comparisons we use a directed lower bound
\begin{equation}\label{app:local:cr}
 c_r=(1-\kappa)c_K>0
\end{equation}
from this intrinsic lemma. All residual lower bounds used in the
certificates follow from this lemma.
The scalar enclosures cover the \emph{whole} reciprocal-$N$ interval,
not merely the local endpoint. They are produced by
\eqref{app:local:intrinsic-scalars}--\eqref{app:local:intrinsic-eigen}
with the root enclosure and contractor described next.

Write $x=1/N$, $p_a=1/4+\delta_ax$, $h_0=\alpha-1/2$,
$\lambda=\alpha-\zeta$, and $h=h_0-\zeta$.
Polynomial division of the defining candidate equation gives exactly
\begin{equation}\label{app:local:root-contractor}
 \zeta=x^2
 \frac{3+1/h-1/(4h^2)}{2+1/(4h_0h)}
                 \sum_a\frac{\delta_a^2}{h-2\delta_ax}.
\end{equation}
For verification,
$\sum_a(h-2\delta_ax)^{-1}
=4/h+(4x^2/h^2)\sum_a\delta_a^2/(h-2\delta_ax)$;
the uniform equation is $2h_0+1-1/(4h_0)=0$.
On any closed $x$ interval start with
$0\le\zeta\le(2/3)Cx_{\rm hi}^2$ and intersect this interval three times
with the outward interval image of \eqref{app:local:root-contractor}.
The true root remains inside each intersection; convergence is not
assumed. Positivity of every denominator is checked before division.
For the actual candidate the discriminant is nonnegative. Its enclosing
interval may therefore be intersected with $[0,\infty)$; a negative lower
endpoint is replaced by zero, while a negative upper endpoint fails the check.
Taking the minimum over all intervals, signs and four vertex types gives
$c_K$. All intervals meeting a requested endpoint are retained whole,
including any interval crossing it.

\subsection{The normal component and its finite budget}
\label{app:local:normal}

The candidate diagonal identity is
$p_a\|z_a\|^2=1/4+(p_a-1/4)/b$ with
$b=3(2\lambda-1)>109/100$. It implies
\[
 \tfrac14\sum_a\|f_a-4p_az_a\|^2
 =\sum_ap_a\|z_a-f_a\|^2+4\sigma^2/b .
\]
Set
\[
 b_{\rm ph}=\sqrt{(B+400/109)\sigma^2}.
\]
For a physical horizontal mean $H$ this bounds its reference constraint
matrix by $b_{\rm ph}\|H\|_\nu$; in particular its reference normal
projection has norm at most $\sqrt2b_{\rm ph}\|H\|_\nu$.
The following lemma uses these bounds in the original frame variables.

\begin{lemma}[Positive normal coupling]\label{app:local:PSD-normal}
Define
\[
 t=\min\{\sqrt2\eta,\sqrt{2B\sigma^2/\omega_p}\},\qquad
 a=\min\{\omega_p^{-1/2},1+t\},
\]
and
\begin{align}
 J_*&=\frac{(1/2+b_{\rm ph})/\sqrt3+(7/8)t+(E_{\mathrm{Hess}}/2)a}{\sqrt{\omega_p}},
 \notag\\
 L_*'&=\frac1{2\sqrt3}+\frac74t+\frac78t^2+(E_{\mathrm{Hess}}/2)a^2,
 \notag\\
 H_{\rm amb}&=(7/4+E_{\mathrm{Hess}})/(2\omega_p),\qquad
 J=\min(J_*,H_{\rm amb}),\quad L=\min(L_*',H_{\rm amb}).
 \label{app:local:JL}
\end{align}
Then
$|Q(H,R_n)|\le NJ\|H\|_F\|R_n\|_F$ and
$Q(R_n)\le NL\|R_n\|_F^2$.
\end{lemma}
\begin{proof}
Write $C_n=cI+D$ with $\tr D=0$ and set
$r=\|R_n\|_F/\sqrt N$. Tightness gives
$r^2=\tr(C_n^2)/2=\|C_nF\|_\nu^2$.
Because $C_n\succeq0$ in dimension two,
$\|DF\|_\nu^2=\tr(D^2)/2\le r^2/2$.
This half-energy estimate would fail for an indefinite normal.
Also $\|C_n\|\le\sqrt2r$, so for $E=C_n(Z-F)$,
$\|E\|_\nu\le tr$, and $\|R_n\|_\nu\le ar$.
Both the point and aggregate alternatives in $t$ refer to this same $E$.

By the reference tensor and the physical normal projection,
\[
 |q_0(H,C_nF)|
 \le\frac{1/2+b_{\rm ph}}{\sqrt3}\|H\|_\nu r .
\]
The error $q_0(H,E)$ is at most $(7/8)t\|H\|_\nu r$.
Add \eqref{app:local:ambient-perturbation} and use
$\|H\|_\nu\le\|H\|_F/\sqrt{N\omega_p}$ to obtain $J_*$.
For the pure normal, the scalar part is null and the traceless part
contributes at most $r^2/(2\sqrt3)$ to $q_0(C_nF)$.
The two terms involving $E$ are bounded by
$(7t/4+7t^2/8)r^2$, and the perturbation contributes
$E_{\mathrm{Hess}} a^2r^2/2$. This gives $L_*'$.
The ambient operator bound independently gives $H_{\rm amb}$ for each.
\end{proof}

Let $c_m>0$ be the accepted mean lower bound and $c_r>0$ the intrinsic
residual lower bound. Put
$h=\|H\|_F/d$, $r=\|U\|_F/d$, $z=\|R_n\|_F/d$.
Although $\Delta=\Xi-R_n$, an absolute mixed upper bound gives
\begin{equation}\label{app:local:normal-optimization}
 \frac{Q(\Delta)}{Nd^2}
 \le-c_mh^2-c_rr^2+2Jzh+Lz^2,\qquad
 h^2+r^2+z^2=1,\quad 0\le z\le\rho_n .
\end{equation}
If $d\le d_c/\sqrt N$, take $\rho_n^2=d_c^2/(2N^2)\le1/2$.
The exact supremum of the right-hand side is
\begin{equation}\label{app:local:normal-max}
 \begin{cases}
 -c_r+(c_r+L+J^2/(c_m-c_r))\rho_n^2,
 &c_m>c_r,\ J^2\rho_n^2\le(c_m-c_r)^2(1-\rho_n^2),\\
 -c_m+(c_m+L)\rho_n^2+2J\sqrt{\rho_n^2(1-\rho_n^2)},
 &\text{otherwise}.
 \end{cases}
\end{equation}
To prove this, fix $z$ and maximize
$-(c_m-c_r)h^2+2Jzh$ on $[0,\sqrt{1-z^2}]$.
The interior value and the boundary value are nondecreasing in $z$
for $z\le1/\sqrt2$, and agree at the transition.
The boundary case includes $c_m\le c_r$, so no division by a nonpositive
gap occurs. Denote the negative of \eqref{app:local:normal-max} by
$g_{\rm normal}$. All branch choices are made with the conservative
rational coefficients actually supplied to the certificate.

\subsection{Angular transport and the shared radial--tangent budget}
\label{app:local:angular}

The columnwise tangent used in the Taylor remainder is different from
the horizontal array $\Xi$ used in the Hessian. Define
\[
 D_r=\frac{(4/11)e}{1-4e},\qquad r_-=\sqrt{1-D_r},\quad
 r_+=\sqrt{1+D_r},\qquad d_c=\sqrt{A_R}+\sqrt{BC}.
\]
The candidate radial identity in Section~\ref{app:candidate} gives
$r_-\le\|z_a\|\le r_+$. Closest alignment and the triangle inequality
give $d\le d_c/\sqrt N$.
In the preliminary common alignment, the exact identity
\[
 \|rv-f\|^2=(r-1)^2+r\,d_{\rm ph}(v,f)^2
\]
for unit $v,f$ in closest phases gives candidate angular bounds
$BC/(Nr_-)$ in sum and $C_{\mathrm{pt}}C/(N^2r_-)$ pointwise. Put
\[
 B_{\rm ang}=\sqrt{L_U\bar E}+\sqrt{BC/r_-},\qquad
 \chi=\sqrt{\frac{(1+s)r_+}{(1-s)r_-}},\qquad
 D_U=\frac23-\frac{L_*}{N}>0.
\]
For fixed common unitary, minimizing individual phases gives
\[
 \|y_i-z_i\|^2=(\|y_i\|-\|z_i\|)^2+
 \|y_i\|\|z_i\|d_{\rm ph}(u_i,z_i/\|z_i\|)^2.
\]
Therefore the actual closest unitary minimizes the angular sum with
the positive weights $\|y_i\|\|z_i\|$. Comparison to the preliminary
alignment yields
\begin{align}
 \sum_i\|y_i\|\|z_i\|d_{\rm ph}(u_i,z_i/\|z_i\|)^2
 &\le (1+s)r_+B_{\rm ang}^2/N,\label{app:local:weighted-angle}\\
 \sum_i d_{\rm ph}(u_i,z_i/\|z_i\|)^2
 &\le\chi^2B_{\rm ang}^2/N. \nonumber
\end{align}
No unweighted angular optimality is used.

For completeness, the rotation estimate needed to transfer the pointwise
bound follows from the tetrahedral second moment. Modulo scalar phase
write $U=\cos\theta I+i\sin\theta(n\cdot\sigma)$ and put $t=|\sin\theta|$.
If $v_a$ are the reference Bloch vectors, then
$|f_a^*Uf_a|^2=1-t^2(1-(n\cdot v_a)^2)$ and
$d_{\rm ph}(Uf_a,f_a)^2\ge t^2(1-(n\cdot v_a)^2)$.
Since $\sum v_av_a^T=4I/3$,
\[
 \sum_i d_{\rm ph}(Uf_{a(i)},f_{a(i)})^2\ge ND_Ut^2.
\]
The phase-distance triangle inequality and
\eqref{app:local:weighted-angle} bound its left side by
$(1+\chi)^2B_{\rm ang}^2/N$. If
\begin{equation}\label{app:local:rotation-guard}
 \frac{(1+\chi)^2B_{\rm ang}^2}{N^2D_U}
 <1-(255/256)^2,
\end{equation}
every $|f_a^*Uf_a|>255/256$, and
\[
 d_{\rm ph}(Uf_a,f_a)^2=
 \frac{2(1-|f_a^*Uf_a|^2)}{1+|f_a^*Uf_a|}
 \le\frac{512}{511}t^2.
\]
Thus one may use
\[
 \mathrm{Rot}_{\rm ang}=
 \sqrt{\frac{512(1+\chi)^2B_{\rm ang}^2}{511D_U}}.
\]
An independent physical-distance triangle gives
$\sum_i d_{\rm ph}(Uf_a,f_a)^2\le4d_c^2/N$ and hence
$\mathrm{Rot}_{\rm phys}=\sqrt{(512/511)4d_c^2/D_U}$,
provided $4d_c^2/(N^2D_U)<1-(255/256)^2$.
Use the smaller of these two bounds only when each selected bound has
its own stated guard; denote the result by $\mathrm{Rot}$.

Write, in the actual individual phases,
\begin{equation}\label{app:local:column-split}
 y_i=(1+t_i)z_i+\xi_i=r_i z_i+\xi_i,\qquad
 z_i^*\xi_i=0,\quad t_i\in\mathbb R.
\end{equation}
Let $X^2=\sum_i\|\xi_i\|^2$ and let $T$ bound $\max_i\|\xi_i\|$.
The following simultaneous choices are valid:
\begin{align}
 q_t&=\frac1N\min\{d_c,(1+s)\sqrt{r_+/r_-}B_{\rm ang}\},
 & X/\sqrt N&\le q_t,\label{app:local:tangent-bounds}\\
 \epsilon&=\sqrt{P_R/N}+(\mathrm{Rot}_{\rm phys}+\sqrt{C_{\mathrm{pt}}C})/N,
 & q_d&=d_c/N,\nonumber\\
 T&=\min\left\{\epsilon,(1+s)\left[
 \sqrt{L_P\bar E/N}+\frac{\mathrm{Rot}}N+
 \frac{\sqrt{C_{\mathrm{pt}}C}}{N\sqrt{r_-}}\right]\right\}.&&\nonumber
\end{align}
Indeed $1-|u^*v|^2\le d_{\rm ph}(u,v)^2$ for unit vectors.
Multiplying \eqref{app:local:weighted-angle} by $(1+s)/r_-$ gives
$X^2\le(1+s)^2(r_+/r_-)B_{\rm ang}^2/N$.
The pointwise estimate follows from the same triangle inequality with
the rotation bound, while the physical triangle gives $\epsilon$.
An angular displacement has therefore not been charged a first-order
radial point error.

The certificate checks $T<1-s$. Since $z_i^*y_i\ge0$ in closest phases,
\begin{equation}\label{app:local:radial-interval}
 0<\rho=\frac{\sqrt{(1-s)^2-T^2}}{r_+}\le r_i
 \le R=\frac{1+s}{r_-}.
\end{equation}
Orthogonality in \eqref{app:local:column-split} gives the exact identity
\begin{equation}\label{app:local:shared-budget}
 d^2=\sum_i t_i^2\|z_i\|^2+X^2,\qquad
 T_r^2:=\sum_i t_i^2\le(d^2-X^2)/r_-^2.
\end{equation}
This budget is distinct from the global radial, moment and angular
deficits. Equations \eqref{app:local:tangent-bounds} are the explicit
transfer between them. For $d>0$ put $\vartheta=X/d\in[0,1]$.

\subsection{The pair remainder and physical fiber bounds}
\label{app:local:pair-remainder}

For an ordered pair put $s_{ij}=|z_i^*Kz_j|$, and rotate its candidate
overlap to the positive real axis. In this phase set
\[
 \alpha_{ij}=\overline S_{ij}z_i^*K\xi_j,\quad
 \beta_{ij}=\overline S_{ij}\xi_i^*Kz_j,\quad
 \gamma_{ij}=\overline S_{ij}\xi_i^*K\xi_j,
 \qquad (a_{ij},b_{ij},c_{ij})=\Ima(\alpha_{ij},\beta_{ij},\gamma_{ij}).
\]
For this paragraph abbreviate $s=s_{ij}$, $p=r_ir_j$ and
$L=r_i\alpha_{ij}+r_j\beta_{ij}$. The complete Taylor polynomial
through total degree two in $(t,\xi)$ is
\[
 P_2=ps+\Rea(L+\gamma_{ij})+(a_{ij}+b_{ij})^2/(2s).
\]
With $R_s(h)=|s+h|-s-\Rea h-(\Ima h)^2/(2s)$, the exact remainder is
\begin{align}
 |ps+L+\gamma_{ij}|-P_2={}&
 \frac{r_i-r_j}{2s}\left(\frac{a_{ij}^2}{r_j}-\frac{b_{ij}^2}{r_i}\right)
 +\frac{(r_i a_{ij}+r_j b_{ij})c_{ij}}{ps}
 +\frac{c_{ij}^2}{2ps}\nonumber\\
 &+pR_s\left(\frac{\alpha_{ij}}{r_j}+
              \frac{\beta_{ij}}{r_i}+\frac{\gamma_{ij}}p\right).
 \label{app:local:exact-pair}
\end{align}
To verify it, expand the square of
$(r_i a_{ij}+r_j b_{ij}+c_{ij})$ divided by $2ps$;
the $a_{ij}b_{ij}$ term cancels exactly. In particular the pure radial
dependence is exactly quadratic. All within-cluster and diagonal pairs
remain in this identity.

For types $\mathrm c$ (different clusters) and $\mathrm s$ (same cluster),
choose positive floors $h_t$ and upper bounds $m_t,g_t$ such that
\begin{equation}\label{app:local:type-bounds}
 s_{ij}\ge h_t,\quad |\alpha_{ij}|\le m_t\|\xi_j\|,\quad
 |\beta_{ij}|\le m_t\|\xi_i\|,\quad
 |\gamma_{ij}|\le g_t\|\xi_i\|\|\xi_j\|.
\end{equation}
We use $h_{\mathrm c}=h_c$ for the certified cross floor.
Besides the certified candidate overlap floor, valid generic choices are
\begin{gather}
 h_E=15/26-2\eta-\eta^2>0,\qquad
 h_{\mathrm s}=(1-\kappa)r_-^2,\nonumber\\
 m_{\mathrm c}=\sqrt{r_+^2-h_E^2/r_+^2}+\kappa r_+,\qquad
 m_{\mathrm s}=\kappa r_+,\label{app:local:generic-fiber}\\
 g_{\mathrm c}=\kappa+\min\{1,(1/\sqrt3+2\eta+\eta^2)/r_-^2\},\qquad
 g_{\mathrm s}=1+\kappa.\nonumber
\end{gather}
The displayed choice of $h_{\mathrm s}$ may be replaced by another
independently proved floor; in all cases require
$0<h_{\mathrm s}\le\min_a s_{aa}$.
The projected-linear bound follows from the Euclidean cross floor and
the component of $K-I$. For the quadratic bound, define the antiunitary
$J(v_1,v_2)=(-\overline v_2,\overline v_1)$. Each
$\xi_i=c_iJ(z_i/\|z_i\|)$, so
\[
 |\xi_i^*\xi_j|=
 \frac{|z_i^*z_j|}{\|z_i\|\|z_j\|}\|\xi_i\|\|\xi_j\|.
\]
Adding the $K-I$ contribution proves $g_{\mathrm c},g_{\mathrm s}$.
In particular this is a physical complex-line calculation, valid for
every within-cluster residual.

\subsection{Row estimates coupled to intrinsic residual curvature}
\label{app:local:rows}

Write $e_i=\|\xi_i\|$. Retain the separate pointwise coefficient
\[
 M=\max\{m_{\mathrm c}^2/h_{\mathrm c},m_{\mathrm s}^2/h_{\mathrm s}\}.
\]
The following row bounds keep the complete imaginary parts in
\eqref{app:local:exact-pair}:
\[
 \sum_i a_{ij}^2/s_{ij}\le NH e_j^2,\qquad
 \sum_{a(i)\ne a(j)}a_{ij}^2/s_{ij}\le NH_{\mathrm c}e_j^2.
\]
To obtain an initial bound, put
\begin{gather}
 d_g=2\eta+\eta^2+\kappa r_+^2,\qquad
 \delta=\eta+\kappa r_++\eta/\sqrt{r_-}
       +\sqrt{2/3}\,d_g/\sqrt{h_{\mathrm c}/\sqrt3},\nonumber\\
 H_{\rm iso}=\frac{(\sqrt{p_+}+\sqrt{1-p_-}\,\delta)^2}{h_{\mathrm c}},
 \label{app:local:row-initial}\\
 H=\min\{M,H_{\rm iso}+p_+m_{\mathrm s}^2/h_{\mathrm s}\},\qquad
 H_{\mathrm c}=\min\{H_{\rm iso},(1-p_-)m_{\mathrm c}^2/h_{\mathrm c}\}.
 \nonumber
\end{gather}
The subscript on $H_{\rm iso}$ denotes the reference isotropy estimate.
For a reference center $f_b$, the three projected cross forms have
amplitude $\sqrt{2/3}\|\xi\|$ and phases separated by $2\pi/3$.
Their squared imaginary parts sum to $\|\xi\|^2$, so their weighted sum
is at most $p_+\|\xi\|^2$. Transport $\xi=cJ(z_b/\|z_b\|)$ to
$cJf_b$; its change is at most $\eta\|\xi\|/\sqrt{r_-}$.
For complex nonzero $v,v_0$,
\[
 \left|\frac v{|v|}-\frac{v_0}{|v_0|}\right|^2
 =\frac{|v-v_0|^2-(|v|-|v_0|)^2}{|v||v_0|}.
\]
The overlap phase changes by at most $d_g/\sqrt{h_{\mathrm c}/\sqrt3}$.
These bounds give the displayed $\delta$ for the full phase-aligned
linear form. The weighted Euclidean triangle inequality proves
$H_{\rm iso}$; same-cluster rows and pointwise $M$ give the other terms.

There is a second, sharper row bound directly from intrinsic curvature.
Let $c_{\rm intr}$ be the physical pointwise coefficient established in
the residual argument above, before imposing any zero-sum condition:
\begin{equation}\label{app:local:pointwise-residual}
 \sum_b p_b\frac{[\Ima(\overline S_{ba}z_b^*Kh)]^2}{|z_b^*Kz_a|}
       -\lambda h^*Kh\le-c_{\rm intr}\|h\|^2,\qquad
 \Ima(z_a^*h)=0.
\end{equation}
Every columnwise $\xi_j$ is eligible since $z_j^*\xi_j=0$. Therefore,
for any certified upper $\lambda_+$,
\begin{equation}\label{app:local:spectral-row}
 H_{\rm spec}=\lambda_+(1+\kappa)-c_{\rm intr},\qquad
 H\leftarrow\min(H,H_{\rm spec}),\quad
 H_{\mathrm c}\leftarrow\min(H_{\mathrm c},H_{\rm spec}).
\end{equation}
The certificate checks $H_{\rm spec}\ge0$. Every summand in the row
is nonnegative, which justifies the cross-only bound. The coefficient
$M$ is not changed: a row bound does not imply a pointwise pair bound.
Nor is a mean curvature coefficient substituted for $c_{\rm intr}$.
Using \eqref{app:local:pointwise-residual} both here and in quadratic
descent is valid because both are pointwise consequences of the same
inequality; neither spends a separate portion of a deficit.

Hermitian symmetry gives $b_{ij}=-a_{ji}$, $c_{ji}=-c_{ij}$.
The full ordered radial term in \eqref{app:local:exact-pair} is consequently
\[
 \sum_{i,j}(t_i-t_j)a_{ij}^2/(r_js_{ij}).
\]
Put $W_i=\sum_j a_{ij}^2/s_{ij}$. Then $0\le W_i\le MX^2$,
$\sum_iW_i\le NHX^2$, hence $\sum_iW_i^2\le NMHX^4$.
Cauchy--Schwarz, followed by \eqref{app:local:shared-budget}, bounds this
term divided by $Nd^2$ by
\begin{equation}\label{app:local:radial-row}
 B_{\rm rad}\vartheta\sqrt{1-\vartheta^2},\qquad
 B_{\rm rad}=\frac{HT+\sqrt{MH}\,q_t}{\rho r_-}.
\end{equation}
Indeed the two sums with $|t_i|$ and $|t_j|$ are at most
$\sqrt{NMH}\,T_rX^2$ and $NHT_rTX$, respectively, before division by
$\rho$.

The full ordered linear--quadratic term is
$2\sum_{i,j}a_{ij}c_{ij}/(r_js_{ij})$. If
$E_b=\sum_{a(j)=b}e_j^2$, its cross part, by weighted Cauchy--Schwarz,
is at most
\[
 \frac{2g_{\mathrm c}}\rho\sqrt{NH_{\mathrm c}/h_{\mathrm c}}
 \sum_b E_b\sqrt{X^2-E_b}
 \le\frac{\sqrt3g_{\mathrm c}}\rho
        \sqrt{NH_{\mathrm c}/h_{\mathrm c}}X^3.
\]
The last inequality follows from
$\sum E_b^2\ge X^4/4$. The same-cluster part is at most
$2\sqrt{Np_+}m_{\mathrm s}g_{\mathrm s}X^3/(\rho h_{\mathrm s})$.
Set
\begin{gather}
 L_{\rm row}=\frac{\sqrt3}2g_{\mathrm c}\sqrt{H_{\mathrm c}/h_{\mathrm c}}
        +\sqrt{p_+}m_{\mathrm s}g_{\mathrm s}/h_{\mathrm s},\qquad
 J_{\max}=\max_t g_t^2/(2h_t),\nonumber\\
 A_{\rm mixed}=2L_{\rm mixed}q_t/\rho+J_{\max}q_t^2/\rho^2,
 \label{app:local:mixed-row}
\end{gather}
where $L_{\rm mixed}$ can be $L_{\rm row}$ or its minimum with the
independent weighted type bound described next. The second term in
$A_{\rm mixed}$ retains the full quadratic--quadratic term $c_{ij}^2$.
The same-cluster term in $L_{\rm row}$ is also retained.

\subsection{Complete scalar bounds and four-cluster summation}
\label{app:local:scalar-tail}

The elementary modulus identities below retain all orders of the
remainder; no assertion that its argument has modulus less than its
base is required. For $h\ne0$ and $s>0$ put
\[
 z=|h|/s,\qquad t=(1-|1+h/s|)/z,\qquad A_t=1-t^2.
\]
Reverse triangle gives $-1\le t\le1$. Squaring the defining identity
gives
\[
 \frac{\Rea h}{s}=-zt-z^2A_t/2,\qquad
 \frac{(\Ima h)^2}{s^2}=z^2A_t-z^3tA_t-z^4A_t^2/4.
\]
Consequently
\begin{equation}\label{app:local:scalar-identity}
 \frac{R_s(h)}s=z^3 f(t)+z^4 g(t),\qquad
 f(t)=\frac{t(1-t^2)}2,\quad g(t)=\frac{(1-t^2)^2}8.
\end{equation}
Replacing $t$ by $|t|$ only increases the right side. On $[0,1]$ put
$a_0=1/\sqrt3$. Then
\[
 f(a_0)-f(t)=\frac{(t-a_0)^2(t+2a_0)}2
 \ge a_0(t-a_0)^2,\qquad
 g(a_0)=1/18,\qquad |g'(t)|\le1/(3\sqrt3).
\]
The mean-value inequality followed by completing a square proves
\begin{equation}\label{app:local:unsigned-scalar}
 R_s(h)\le \frac{|h|^3}{3\sqrt3s^2}
       +\left(\frac1{18}+\frac Z{36\sqrt3}\right)\frac{|h|^4}{s^3},
 \qquad |h|/s\le Z.
\end{equation}
At $h=0$ the assertion is immediate.

Define the fixed mass coefficient $\zeta=(3/2)\sqrt{p_+}<1$.
For any nonnegative frozen pair of coefficients
$F_{\mathrm c}\ge F_{\mathrm s}$ set
\begin{equation}\label{app:local:mass-weights}
 F_{\rm eff}=(1-p_-)F_{\mathrm c}+p_-F_{\mathrm s},\quad
 F_{\rm mix}=\zeta F_{\mathrm c}+(1-\zeta)F_{\mathrm s},\quad
 F_{\rm pair}=(3F_{\mathrm c}+F_{\mathrm s})/4 .
\end{equation}
Here and below, ordering is checked for the chosen majorants, not
inferred from an ordering of the actual geometric quantities.
The ordered cubic cross sum satisfies
\[
 \sum_{a(i)\ne a(j)}e_i^2e_j\le\zeta\sqrt N X^3,\qquad
 \sum_{i,j}e_i^2e_j\le\sqrt N X^3.
\]
Indeed, if $x_a=\sqrt{\sum_{a(i)=a}e_i^2}/X$, then $\sum x_a^2=1$.
The first sum is at most
$\sqrt{Np_+}X^3(\sum x_a-\sum x_a^3)$.
The inequalities $\sum x_a\le2$ and
$\sum x_a^3\ge(\sum x_a)/4$ prove the claim. The latter follows from
$4\sum x_a^3\ge(\sum x_a^2)(\sum x_a)$; expanding gives a sum of
$(x_a-x_b)(x_a^2-x_b^2)\ge0$.
Also $\sum_{a(i)\ne a(j)}e_i^2e_j^2\le3X^4/4$, since
$\sum_aE_a^2\ge X^4/4$. Pure powers are weighted by the number of
cross partners, at most $N(1-p_-)$.
These proofs use arbitrary column energies and do not assume
constant perturbations in a cluster.

In particular the independent type bound used in
\eqref{app:local:mixed-row} is $(L_t)_{\rm mix}$ with
$L_t=m_tg_t/h_t$; take its minimum with $L_{\rm row}$ when
$L_{\mathrm c}\ge L_{\mathrm s}$.
For the scalar term of \eqref{app:local:exact-pair}, put
\begin{equation}\label{app:local:unsigned-types}
 v_t=\frac{g_tT}{2\rho},\quad D_t=m_t+v_t,\quad
 Z_t=\frac{2D_tT}{\rho h_t},\quad
 C_t=\frac{D_t^3}{3\sqrt3h_t^2},\quad
 B_t=\frac{D_t^4}{h_t^3}\left(\frac1{18}+\frac{Z_t}{36\sqrt3}\right).
\end{equation}
The argument of $R_s$ has modulus at most
$D_t(e_i/r_i+e_j/r_j)$, since
$e_ie_j/(r_ir_j)\le T(e_i/r_i+e_j/r_j)/(2\rho)$.
Keep $r_ir_j$ outside the scalar remainder until expanding the powers:
\begin{align*}
 r_ir_j(e_i/r_i+e_j/r_j)^3
 &=\frac{r_j}{r_i^2}e_i^3+\frac{r_i}{r_j^2}e_j^3
       +3e_i^2e_j/r_i+3e_ie_j^2/r_j,\\
 r_ir_j(e_i/r_i+e_j/r_j)^4
 &=\frac{r_j}{r_i^3}e_i^4+\frac{r_i}{r_j^3}e_j^4
       +4e_i^3e_j/r_i^2+4e_ie_j^3/r_j^2
       +6e_i^2e_j^2/(r_ir_j).
\end{align*}
Using \eqref{app:local:mass-weights} gives the full unsigned coefficients
\begin{align}
 A_3^{\,u}&=2RC_{\rm eff}T/\rho^2+6C_{\rm mix}q_t/\rho,\nonumber\\
 A_4^{\,u}&=2RB_{\rm eff}T^2/\rho^3
             +8B_{\rm mix}Tq_t/\rho^2+6B_{\rm pair}q_t^2/\rho^2,
 \label{app:local:unsigned-loss}\\
 A^{\,u}&=A_3^{\,u}+A_{\rm mixed}+A_4^{\,u},\qquad B^{\,u}=B_{\rm rad}.
 \nonumber
\end{align}
The necessary $C$ and $B$ orderings are part of the certificate guards.

\begin{lemma}[Shared-budget remainder]\label{app:local:loss}
Suppose one complete scalar comparison supplies coefficients $A,B\ge0$
such that the ordered remainder in \eqref{app:local:exact-pair}, divided
by $Nd^2$, is at most
$A\vartheta^2+B\vartheta\sqrt{1-\vartheta^2}$.
Then it is at most
\begin{equation}\label{app:local:loss-formula}
 {\rm Loss}(A,B)=\frac{A+\sqrt{A^2+B^2}}2.
\end{equation}
In particular \eqref{app:local:unsigned-loss} supplies such a comparison.
\end{lemma}
\begin{proof}
The maximum is the larger eigenvalue of
$\left(\begin{smallmatrix}A&B/2\\B/2&0\end{smallmatrix}\right)$.
Equations \eqref{app:local:radial-row},
\eqref{app:local:mixed-row} and the preceding power expansions establish
the unsigned assertion with the same $\vartheta$ in every term.
\end{proof}

\subsection{A signed remainder bound}
\label{app:local:signed}

Let $P_s(h)=-(\Rea h)(\Ima h)^2/(2s^2)$ be the scalar cubic.
Substitution in \eqref{app:local:scalar-identity} yields the exact identity
\begin{equation}\label{app:local:compact-identity}
 \frac{R_s(h)-P_s(h)}s=z^4H_z(t),\qquad
 H_z(t)=\frac{(1-t^2)(5t^2-1)}8
       +\frac{3z}{8}t(1-t^2)^2+\frac{z^2}{16}(1-t^2)^3 .
\end{equation}
This identity includes every order in $h$. Replacing $t$ by $|t|$
raises its right side, and $H_z(t)$ is nondecreasing in $z$ for
$0\le t\le1$. More generally choose a fixed sign weight
$\omega\in[0,1]$ and a rational $c_3\ge1/(3\sqrt3)$. For $Z>0$ let
\[
 W_{Z,\omega}(t)=
 \frac{(1-\omega)(f(t)-c_3)}Z+(1-\omega)g(t)+\omega H_Z(t).
\]
If $Q_\omega(Z)\ge\max_{0\le t\le1}W_{Z,\omega}(t)$, then
\begin{equation}\label{app:local:fractional-scalar}
 R_s(h)\le\omega P_s(h)+(1-\omega)c_3|h|^3/s^2
                       +Q_\omega(Z)|h|^4/s^3,\quad |h|/s\le Z.
\end{equation}
To verify it, subtract the first two terms, divide by $sz^4$, and use
\eqref{app:local:compact-identity}. The result is
$(1-\omega)(f(t)-c_3)/z+(1-\omega)g(t)+\omega H_z(t)$.
Since $f(t)\le c_3$, replacing $0<z\le Z$ by $Z$ increases it.
The zero case is direct. No condition $Z<1$ enters this proof.

The continuum maximum is bounded by rational Bernstein coefficients.
The ascending power coefficients of $H_Z$ are
\[
 \left(\frac{-2+Z^2}{16},\frac{3Z}{8},
 \frac{12-3Z^2}{16},-\frac{3Z}{4},
 \frac{-10+3Z^2}{16},\frac{3Z}{8},-\frac{Z^2}{16}\right).
\]
Add the displayed $f,g,c_3$ terms to obtain $W_{Z,\omega}$.
On a closed interval $[a,b]\subset[0,1]$, substitute
$t=a+(b-a)v$ and call the resulting power coefficients $c_j$.
Its degree-six Bernstein coefficients are
\[
 \beta_i=\sum_{j=0}^i c_j\frac{\binom ij}{\binom6j},\qquad 0\le i\le6.
\]
The identity
$W_{Z,\omega}(a+(b-a)v)=\sum_i\beta_i\binom6iv^i(1-v)^{6-i}$
and the nonnegative partition of unity bound the whole interval by
$\max_i\beta_i$. The certificate uses the $32$ closed equal intervals.
It checks $Q_\omega(Z)\ge0$ before substituting the nonnegative power
majorants in the ordered sums.

We next explain the pure-row cubic input; its sharpened version is
proved in the following subsection. For a fixed column $j$ choose a
unit vector $v_j\perp z_j$ and set
\[
 \ell_{ij}=\overline S_{ij}z_i^*Kv_j,\qquad
 q_{ij}=|\ell_{ij}|^2/s_{ij}^2 .
\]
Stationarity, with its full stress term, is
\begin{equation}\label{app:local:cubic-stationarity}
 \frac1N\sum_i\ell_{ij}=\lambda z_j^*Kv_j,\qquad
 |\lambda z_j^*Kv_j|\le\lambda_+\min\{\kappa r_+,m_{\mathrm s}\}.
\end{equation}
Its right side is generally nonzero. The Fourier identity
\[
 P_s(\ell u)=
 \Rea\{\ell^3u^3-|\ell|^2\ell|u|^2u\}/(8s^2)
\]
shows that a pure row bound follows from its first and third harmonics.
Given $q_{\mathrm c}\in[q_-,q_+]$, $q_{\mathrm s}\in[0,q_{\mathrm s,+}]$
and any fixed $c\ge0$, define
\begin{align}
 F_1(c)={}&c\lambda_+\min\{\kappa r_+,m_{\mathrm s}\}
 +(1-p_-)\max\{|q_--c|,|q_+-c|\}m_{\mathrm c}\nonumber\\
 &+p_+\max\{c,|q_{\mathrm s,+}-c|\}m_{\mathrm s},\nonumber\\
 F_3={}&(1-p_-)m_{\mathrm c}^3/h_{\mathrm c}^2
                  +p_+m_{\mathrm s}^3/h_{\mathrm s}^2,\qquad
 P_{\rm row}=(F_1(c)+F_3)/8 .\label{app:local:generic-cubic-row}
\end{align}
Subtract $c$ times \eqref{app:local:cubic-stationarity} from the first
harmonic and apply the triangle inequality; the third harmonic uses
its full typed triangle bound. Thus
\begin{equation}\label{app:local:pure-row}
 N^{-1}\sum_iP_{s_{ij}}(\alpha_{ij})\le P_{\rm row}e_j^3.
\end{equation}
Valid generic ratio intervals are
\begin{align*}
 q_-&=\max\left\{0,\frac{(1-\kappa)^2r_-^2}{(1/\sqrt3+d_g)^2}
                              -r_-^{-2}\right\},\\
 q_+&=\min\left\{\frac{m_{\mathrm c}^2}{h_{\mathrm c}^2},
       \frac{(1+\kappa)^2r_+^2}{h_{\mathrm c}^2}-r_+^{-2}\right\},
 \qquad q_{\mathrm s,+}=m_{\mathrm s}^2/h_{\mathrm s}^2 .
\end{align*}
They follow from the exact physical projection identity proved below.
The certificate can intersect these intervals with sharper valid
intervals and choose $c=2$ or their endpoint midpoint. Its chosen $c$
is fixed for the infinite-range argument.

To retain all mixed cubics, the symmetric trilinear form of $P_s$ obeys
$|\mathcal T_s(u,u,v)|\le|u|^2|v|/(3\sqrt3s^2)$.
For $u=x+iy$, the numerator norm is
$\sqrt{y^4+4x^2y^2}/6$, whose maximum at fixed $|u|$ gives this bound.
Expansion of the cubic, or integration of its gradient, gives
\[
 P_s(u+v+w)\le P_s(u)+P_s(v)+
 \frac{(|u|+|v|+|w|)^3-|u|^3-|v|^3}{3\sqrt3s^2}.
\]
All coefficients of the last polynomial are nonnegative.
For each pair type retain $v_t,D_t,Z_t$ from
\eqref{app:local:unsigned-types} and set
\begin{align}
 C_{{\rm base},t}&=c_3m_t^3/h_t^2,\nonumber\\
 C_{{\rm extra},t}&=
 c_3(3m_t^2v_t+3m_tv_t^2+v_t^3)/h_t^2
                      +(1-\omega)C_{{\rm base},t},\nonumber\\
 C_{{\rm full},t}&=c_3(m_t+v_t)^3/h_t^2,\qquad
 B_{{\rm signed},t}=D_t^4Q_\omega(Z_t)/h_t^3.
 \label{app:local:signed-types}
\end{align}
In particular $C_{\rm full}$ does not receive an additional
$(1-\omega)C_{\rm base}$ term. Outward rounding may increase $D_t,Z_t$;
the scalar certificate then uses that rounded radius.

The physical radial factors require a separate estimate. The two pure
ordered cubic terms combine to
$2\sum_jr_j^{-2}\sum_ir_iP_{s_{ij}}(\alpha_{ij})$.
One cannot pull a positive upper bound for $r_i$ out of this signed sum.
Write $r_i=1+t_i$ and apply \eqref{app:local:pure-row} only to the part
with coefficient one. With
\[
 C_{\rm rms}=\sqrt{(1-p_-)C_{{\rm base},{\mathrm c}}^2+
                         p_-C_{{\rm base},{\mathrm s}}^2},
 \qquad B_{\rm weight}=\frac{2\omega C_{\rm rms}Tq_d}{\rho^2r_-},
\]
Cauchy--Schwarz bounds the remaining part by
$B_{\rm weight}\vartheta^2\sqrt{1-\vartheta^2}$ after division by $Nd^2$.
Indeed $\sum_i|t_i|C_{{\rm base},t(i,j)}
\le\sqrt N C_{\rm rms}T_r$ and $\sum_j e_j^3\le TX^2$.
It is safe to replace $\vartheta^2$ by $\vartheta$ here.

Using \eqref{app:local:mass-weights}, the complete signed alternative is
\begin{align}
 A_3^{\,s}&=\frac{2\omega P_{\rm row}T}{\rho^2}
      +\frac{2R C_{{\rm extra},{\rm eff}}T}{\rho^2}
      +\frac{6C_{{\rm full},{\rm mix}}q_t}{\rho},\nonumber\\
 A_4^{\,s}&=\frac{2RB_{{\rm signed},{\rm eff}}T^2}{\rho^3}
      +\frac{8B_{{\rm signed},{\rm mix}}Tq_t}{\rho^2}
      +\frac{6B_{{\rm signed},{\rm pair}}q_t^2}{\rho^2},\nonumber\\
 A^{\,s}&=A_3^{\,s}+A_{\rm mixed}+A_4^{\,s},\qquad
 B^{\,s}=B_{\rm rad}+B_{\rm weight}.
 \label{app:local:signed-loss}
\end{align}
The certificate checks cross/same ordering for $C_{\rm base}$,
$C_{\rm extra}$, $C_{\rm full}$ and $B_{\rm signed}$, and nonnegative
quartic coefficients. Lemma~\ref{app:local:loss} then applies.
Taking the smaller of complete unsigned and complete signed bounds,
including the four fixed choices $\omega=1/4,1/2,3/4,1$, is valid.
Their cubic and scalar-tail pieces are not selected independently.

\subsection{Exact two-pair fibers and row cancellation}
\label{app:local:two-pair}

Here the first index of a fiber denotes its tangent center. Retain
$\mathcal T,d_a,\ell_a$ from \eqref{app:local:intrinsic-scalars}, and put
\begin{gather}
 t_{2a}=d_a^2+\sum_{b\ne a}\ell_a^2\ell_b^2p_ap_b,\qquad
 W_a=\frac{2t_{2a}}{p_a(\tr\mathcal T)^2},\qquad
 R_a=\frac{3\lambda-2+2p_a}{6p_a(2\lambda-1)}.
 \label{app:local:physical-fiber-data}
\end{gather}
Then $W_a=\|Kz_a\|^2$ and $R_a=\|z_a\|^2$.
Indeed $\mathcal T UU^*\mathcal T=\mathcal T^2$ since $UU^*$ projects
onto its range; its diagonal gives the formula for $W_a$.
The formula for $R_a$ is the candidate diagonal identity
$p_aR_a=1/4+(p_a-1/4)/[3(2\lambda-1)]$.
Complex dimension two implies that for every unit $v\perp z_a$,
\begin{equation}\label{app:local:exact-physical-fiber}
 m_{ab}^2:=|z_b^*Kv|^2=W_b-s_{ab}^2/R_a,\qquad
 q_{ab}=m_{ab}^2/s_{ab}^2 .
\end{equation}
This follows by subtracting the projection of $Kz_b$ on
$z_a/\|z_a\|$. In particular the modulus is independent of the choice
of unit vector in this one-dimensional complex fiber, without removing
its phase. For $a=b$ the overlap is $s_{aa}=d_a/(p_a\tr\mathcal T)$,
not the cross formula $\ell_a\ell_b/\tr\mathcal T$.

For the profile $C=1$ put
\begin{equation}\label{app:local:two-pair-parameter}
\begin{gathered}
 x=1/N,\quad t=\sqrt{3-8x^2},\quad
 a=\tfrac14-\tfrac x2,\quad b=\tfrac14+\tfrac x2,\\
 (p_1,p_2,p_3,p_4)=(a,a,b,b),\quad \lambda=(1+t)/4.
\end{gathered}
\end{equation}
Substitution in \eqref{alg:root-equation}, using $t^2+8x^2=3$,
proves the last identity by the uniqueness of the root in $(2/3,\alpha]$.
Reversing the signed count profile permutes the two pairs.
There are six ordered types: $\mathrm{aa},\mathrm{ab},\mathrm{ba},
\mathrm{bb},\mathrm{same}_a,\mathrm{same}_b$.
Here $\mathrm{aa}$ means two different clusters of mass $a$ and must
not be confused with $\mathrm{same}_a$.
The following explicit formulas make their interval verification
independent of a choice of spinor phases:
\begin{align}
 s_{\mathrm{aa}}&=\frac{t(2x+1)}{(2x-1)(8x^2-3)},&
 s_{\mathrm{ab}}&=\frac{-t}{8x^2-3},&
 s_{\mathrm{same}_a}&=
       -\frac{2tx+8x^2-3}{(2x-1)(8x^2-3)}. \label{app:local:sharp-s}
\end{align}
Set $V=4tx+20x^2+8x-3$ and
$J=-32tx^3-12tx^2+6tx+16x^4+56x^3-36x^2-24x+9$.
Then
\begin{align}
 m_{\mathrm{aa}}^2&=
 \frac{-2V}{(2x-1)(44x^2+12x-9)},\nonumber\\
 m_{\mathrm{ab}}^2&=
 \frac{2J}{(2x+1)(8x^2-3)(44x^2+12x-9)},\nonumber\\
 m_{\mathrm{same}_a}^2&=
 \frac{8x^2V}{(2x-1)(8x^2-3)(44x^2+12x-9)}, \label{app:local:sharp-m}\\
 q_{\mathrm{aa}}&=
 \frac{2(2x-1)(8x^2-3)V}{(2x+1)^2(44x^2+12x-9)},\nonumber\\
 q_{\mathrm{ab}}&=
 \frac{-2J}{(2x+1)(44x^2+12x-9)},\nonumber\\
 q_{\mathrm{same}_a}&=
 \frac{8x^2(-32tx^2+104x^3-36x^2-42x+9)}
 {9(2x-1)(2x+1)(44x^2+12x-9)}. \label{app:local:sharp-q}
\end{align}
All formulas for center $b$ are obtained by replacing $x$ by $-x$
and interchanging the labels $a,b$. These identities follow by
substitution in \eqref{app:local:physical-fiber-data} and
\eqref{app:local:exact-physical-fiber}, followed by multiplication by
the denominators and reduction using $t^2=3-8x^2$.
This supplies direct polynomial checks of each identity. The factor
$x^2$ in both same-type formulas is retained before interval evaluation.

For a center of mass $a$, denote the four complex forms in
\eqref{app:local:cubic-stationarity} by $\ell_0$ in its own cluster,
$\ell_1$ in the other $a$ cluster, and $\ell_2,\ell_3$ in the $b$
clusters. Exact stationarity reads
\[
 a\ell_0+a\ell_1+b\ell_2+b\ell_3=\lambda\ell_0.
\]
The two $b$ ratios are the same real number $q_{\mathrm{ab}}$ even
though their complex forms need not have the same phase. Thus the
first harmonic has the exact reduction
\begin{equation}\label{app:local:two-pair-harmonic}
 \sum_jp_jq_j\ell_j
 =a(q_{\mathrm{aa}}-q_{\mathrm{ab}})\ell_1+
 [a q_{\mathrm{same}_a}
                  +q_{\mathrm{ab}}(\lambda-a)]\ell_0.
\end{equation}
There is no assertion that the individual $b$ terms vanish.
Since $\lambda>a$,
\begin{align}
 F_{1a}&=a|q_{\mathrm{aa}}-q_{\mathrm{ab}}|m_{\mathrm{aa}}
    +[a q_{\mathrm{same}_a}+q_{\mathrm{ab}}(\lambda-a)]
                                          m_{\mathrm{same}_a},\nonumber\\
 F_{3a}&=a q_{\mathrm{same}_a}m_{\mathrm{same}_a}
         +a q_{\mathrm{aa}}m_{\mathrm{aa}}
         +2b q_{\mathrm{ab}}m_{\mathrm{ab}},\nonumber\\
 P_{\rm row}^{\,\rm pair}
   &=\max\{F_{1a}+F_{3a},F_{1b}+F_{3b}\}/8
 \label{app:local:two-pair-row}
\end{align}
is a valid bound in \eqref{app:local:pure-row}.
The third harmonic uses $|\ell_j|^3/s_j^2=q_jm_j$ exactly.
The $b$ center uses the exchanged formulas. In particular the
same-cluster stress term survives, and all physical residual vectors
are eligible. One may take the minimum with
\eqref{app:local:generic-cubic-row}.

For clarity, the finite continuum certificate for these formulas uses
the $512$ closed intervals
$[j/(180\cdot512),(j+1)/(180\cdot512)]$, $0\le j<512$.
On each interval $[u,v]$ use
$\sqrt{3-8v^2}\le t\le\sqrt{3-8u^2}$ with directed rational square roots.
Rational interval substitution in
\eqref{app:local:sharp-s}--\eqref{app:local:sharp-q}, for all six types,
checks that every denominator excludes zero. Geometric nonnegativity
permits clipping a lower bound for $m^2,q$ to zero; a negative upper
bound fails the check. Evaluate \eqref{app:local:two-pair-row} on each
same interval, retaining both centers.
The complete recorded bounds on $0\le x\le1/180$ include
\begin{align}
 m_{\mathrm c}^2&\le\frac{674866483913819}{10^{15}},&
 m_{\mathrm s}^2&\le\frac{1730019549}{62500000000000},\nonumber\\
 \frac{236851236625637}{125000000000000}
 &\le q_{\mathrm c}\le\frac{2109213874049103}{10^{15}},&
 q_{\mathrm s}&\le\frac{13969617457}{500000000000000},\nonumber\\
 P_{\rm row}^{\,\rm pair}&\le\frac{160219816186209}{10^{15}}.
 &&\label{app:local:sharp-uniform}
\end{align}
At the local endpoint $182$, every whole interval intersecting
$[0,1/182]$ is retained, including a crossing boundary interval.
The slightly sharper pure-row cap used there is
$32029714476747/200000000000000$.
This is an interval cover, not a maximum over sampled configurations.
Only the pure-row term is replaced in \eqref{app:local:signed-loss};
$B_{\rm weight}$ and every other remainder term are still present.

\subsection{Propagation and equality}
\label{app:local:closure}

We now extend the local comparison from a certified endpoint to all
later integers of the same profile. On each closed localization region,
fix its $U,W,L_U,L_P,\bar E$ and the simultaneous candidate
coefficients $B,C_{\mathrm{pt}}$ supplied by Sections~\ref{app:global} and
\ref{app:candidate}. At an accepted endpoint $n$, put $e_n=m_*/n$ and
$\sigma_n^2=C/n^2$. Let $\kappa_*$ be the outward rational upper enclosure
of $n$ times the explicit stress bound in \eqref{app:candidate:kappa}:
\[
 \kappa_*\ge
 2\sqrt C\sqrt{1/3+4e_n/b_0+4\sigma_n^2/b_0^2}
       +\frac{2ne_n^2}{1/15-e_n},\qquad b_0=109/100.
\]
This coefficient is distinct from the Bloch-Gram coefficient $k_0$ in
Section~\ref{app:candidate}. We retain that accepted enclosure throughout
propagation. At this endpoint, valid physical coefficients are
\begin{equation}\label{app:local:frozen-physical}
 s_n=\sqrt{13U/(4n)},\qquad
 A_R=\frac{13U}4+L_UW\frac{1+s_n}{(1-s_n)^2},\qquad
 P_R=\frac{13U}4+L_PW\frac{1+s_n}{(1-s_n)^2}.
\end{equation}
These follow from the radial/angular identity, with the relative
angular bound $L_Uw/[N(1-s)^2]$ and its pointwise analogue. Their minimum
with an independently certified bound in the original variables is
also valid. The partition and the closed count cover remain fixed, so
the count exclusions apply at every later local endpoint.

At an accepted endpoint, all of the following are checked by outward
rational inequalities: positive masses and overlap floors; the
candidate point domains; every fine phase guard and the separate
closest-gauge rotation guard; the mean-coordinate phase, projection,
and positivity guards; the intrinsic residual lower bound; positive
radial multipliers; both finite-normal branches; the nonnegative
quartic and cross/same ordering guards; and the full remainder margin.
The square root of a positive rational is enclosed by rational numbers
whose squares straddle it, so these checks require no floating-point
eigenvalue or optimizer. Section~\ref{app:certificates} specifies
the interval records, arithmetic and source binding.

\begin{lemma}[Frozen endpoint propagation]\label{app:local:propagation}
Suppose all stated guards and a strict complete local margin hold on
one closed region at endpoint $n$, with intrinsic and exact-fiber
enclosures valid on $0\le1/N\le1/n$. Then the same margin lower bound
holds on that region for every later admissible member of its profile.
\end{lemma}
\begin{proof}
For $N\ge n$, retain the fixed profile deviations $\delta_a$ and use
the actual masses $p_a=1/4+\delta_a/N$. The candidate is reconstructed
at those masses; it is not obtained by scaling a finite-endpoint vector
configuration. Its errors obey
\[
 e=m_*/N,\qquad \eta\le\sqrt{C_{\mathrm{pt}}C}/N,\qquad
 \kappa\le\kappa_*/N,\qquad \tau\le(2/3)C/N^2 .
\]
For the fixed count profile, the stress estimate follows because $N$
times the explicit bound in \eqref{app:candidate:kappa} is nondecreasing in $e=m_*/N$ and
$\sigma^2=C/N^2$; its value for $N\ge n$ is therefore bounded by the
accepted enclosure $\kappa_*$.
Freeze every favorable lower coefficient and every unfavorable upper
coefficient of the accepted endpoint. Here is the justification for
each part of this operation.

First, the global exclusions remain strict and the regional fine
bounds remain valid by Section~\ref{app:global}. In
\eqref{app:local:frozen-physical}, $s_n$ and $(1+s_n)/(1-s_n)^2$
decrease as the endpoint increases. The candidate errors decrease,
$r_-$ increases, and $r_+$ decreases. Thus all angular transport
majorants, $T,q_t,q_d$ and $R$ have their endpoint upper bounds;
$\rho$ has its endpoint lower bound. Both actual-rotation guards remain
valid. No difference of decreasing functions is asserted to be
monotone.

Second, each nonnegative perturbation, phase-defect and projection
majorant in the mass and stress mean estimates decreases when these
candidate errors decrease. Their favorable metric factors increase.
Hence the accepted physical mean lower bound remains valid. The
intrinsic $c_K$ was certified on the entire reciprocal-$N$ domain,
and multiplication by $1-\kappa$ preserves its endpoint lower bound.
The original physical normal estimates similarly retain their fixed
upper $J,L$. The normal radius $d_c/(\sqrt2N)$ shrinks. Therefore
the supremum in \eqref{app:local:normal-optimization} cannot increase:
its coefficient bounds are frozen and its domain only shrinks.
This argument does not keep a particular active maximizing branch.

Third, all generic projected and phase-row majorants have the displayed
endpoint bounds. The exact two-pair inequalities cover the whole
reciprocal domain, without needing monotonicity of an individual
ratio. Freeze $\lambda_+$ and $c_K$ separately in the spectral row
bound; then
$\lambda_+(1+\kappa)-(1-\kappa)c_K$ has its endpoint upper bound.
The cross-only row shares this bound, while the independent pair
coefficient $M$ is retained.

Finally freeze every scalar radius cap $Z_t$, sign weight $\omega$,
Bernstein upper bound $Q_\omega(Z_t)$, harmonic center $c$, and the
separate positive cross/same type coefficients. All later increments
lie in the same certified disks. The nonnegative type polynomials and
the signed physical-weight estimate give the endpoint upper
coefficients in \eqref{app:local:signed-loss}; the unsigned bound
propagates in the same way. The ordering checks concern these fixed
majorants. No monotonicity of a rounded Bernstein output or of a
selected numerical branch is needed.
\end{proof}

The rational certificates in Section~\ref{app:certificates:acceptance}
cover all $7320$ budget cells by strict count exclusions and $228$
accepted closed local regions. On every accepted region, the certified
inequality \eqref{app:certificates:accepted-margin} gives
$g_{\rm normal}-{\rm Loss}>1/10000$.
The detailed coverage totals and rational margins are recorded there.

\begin{proposition}[Uniform strict local comparison]
\label{app:local:strict}
For every nondivisible integer $N\ge180$ and every localized competitor
with its actual closest candidate and gauge,
\begin{equation}\label{app:local:strict-inequality}
 \Gamma_K(Y)\le N^2\lambda-\frac{N}{10000}\|Y-Z\|_F^2 .
\end{equation}
\end{proposition}
\begin{proof}
If $d=0$, stationarity gives equality in
\eqref{app:local:strict-inequality} directly. Otherwise, since
$D\mathcal L(Z)=0$, the second-order Taylor polynomial of
$\mathcal L(Z+\Delta)-\mathcal L(Z)$ is $Q(\Delta)$.
The quadratic estimate is at most $-Ng_{\rm normal}d^2$, and the
ordered Taylor remainder is at most $N\,{\rm Loss}\,d^2$.
Every retained closed region has
$g_{\rm normal}-{\rm Loss}>1/10000$ by the rational checks just stated.
The preceding propagation argument proves the entire subsequent
range. Odd integers begin at $181$; residue-two integers begin at
$182$, so the local endpoint $182$ for $C=1$ leaves no integer gap.
\end{proof}

The comparison applies to every localized kernel, with no restriction
on the number of distinct column directions.
In the ordered global trace chain, equality in the optimal bound
therefore forces $d=0$ by Proposition~\ref{app:local:strict}.
Strict integer balancing then forces the balanced profile, and every
nonzero candidate stress overlap fixes its corresponding matrix entry.
This gives the unitary, individual-phase, permutation and conjugation
classification in the main theorem. The divisible
integers, including $180$, are handled by the exact divisible case of
the global theorem. Hence $180$ is a sufficient threshold.

\section{Exact certificates and the computational proof}\label{app:certificates}
The computer-assisted part of the proof consists of finite rational
inequalities on closed domains. We specify the domain covers, the
acceptance conditions and their propagation to infinitely many integers.
The accompanying reproduction package contains the fixed certificates
and an arithmetic verifier that checks them without repeating the search.

\subsection{Arithmetic and acceptance rules}
All arithmetic used for acceptance is rational. To enclose
$\sqrt{a/b}$, for nonnegative integers $a$ and positive $b$, choose a
positive integer scale $S$ and compute
\[
 q=\left\lfloor\sqrt{\left\lfloor aS^2/b\right\rfloor}\right\rfloor.
\]
The lower endpoint is $q/S$. If $bq^2=aS^2$, the upper endpoint is also
$q/S$; otherwise it is $(q+1)/S$, strictly above $\sqrt{a/b}$.
Interval addition, multiplication,
reciprocal and square roots are rounded outward. Division is accepted
only after its denominator interval excludes zero, and every phase,
overlap and covariance guard is checked before use. Fixed decimal grids
are rational denominator grids, not floating-point error tolerances.
The stored scales and exact endpoints are part of the certificates.

A scalar strict inequality is accepted only when the rational upper
bound is strictly below the rational lower bound on the other side.
We use $1/10000$ as a common certified lower bound for the local margin,
independently of any auxiliary target used to refine an interval cover.
In particular, the continuous residual records marked \texttt{FAIL\_TARGET}
retain rigorous positive lower bounds. The local proof uses these bounds,
regardless of whether they meet the auxiliary target.

\subsection{Global and continuous-parameter certificates}
The anchor potential \eqref{app:global:potential} uses $256$ closed
intervals per budget. With $K_A=441/100$ and $\eta_A=8/25$, its saved
upper maximum is
\[
 \frac{1701784823949001620590499470021431663618000000000000000}
 {1703019875451378709607322015228880928794132391894118139}<1.
\]
The same records provide the strict coverage, disjointness and scatter
inequalities, and the separate count checks for $M=1/2$ and $11/19$.
The fine chain \eqref{app:global:chain} has $9$ recorded stages, each
with $128$ closed budget bins. Its $38$ explicit guards include all
positive denominators and forward implications. Regional fine records
then retain their own $64$-bin bounds and phase choices.

The candidate spectra and overlap floors cover
$0\le x=1/N\le1/180$. The cover is the union of the saved
$[0,1/200]$ domain and the closed strip $[1/200,1/180]$.
The latter has $2048$ overlap boxes and $16245$ residual leaves, with
all four profiles and both signs. The shared endpoint is present in
both domains. A later endpoint $n$ uses whole saved cells meeting
$[0,1/n]$, including the whole of a boundary-crossing cell. Its extra
portion only weakens the bound.

The root enclosures follow from the scalar equation. Put
$h_0=\alpha-1/2$, $\tau=\alpha-\lambda$ and $h=h_0-\tau$.
Exact division of the scalar root equation gives
\begin{equation}\label{app:certificates:root}
 \tau=\frac{x^2W(h)}{D(h)}
       \sum_a\frac{\delta_a^2}{h-2\delta_ax},\qquad
 W(h)=3+1/h-1/(4h^2),\quad D(h)=2+1/(4h_0h).
\end{equation}
Start with $0\le\tau\le(2/3)C x_{\rm hi}^2$ and intersect this interval
with the interval evaluation of \eqref{app:certificates:root} three
times. Existence and uniqueness of the root have already been
proved analytically. Each intersection contains that root because it
satisfies the displayed identity; no contraction factor or convergence
assumption is needed. Every before, right-side and after interval is
stored, and an empty intersection is a failure.

For the balanced residue-two profile, the sharper formula
$\lambda=(1+\sqrt{3-8x^2})/4$ is used for the six ordered projected
fiber types in Section~\ref{app:local}. The exact rational expressions
are enclosed on $512$ closed intervals in $[0,1/180]$. The stationary
row certificate reads these intervals and bounds both possible center
types on the same bins. Sign reversal interchanges the two mass types.
The domain covers justify the physical $m,q$ and $P_{\rm row}$ bounds
throughout the stated range, without a monotonicity assumption.

\subsection{Closed budget covers and local comparisons}
\label{app:certificates:acceptance}
For each of the four count profiles divide the $u,w$ triangle into the
closed cells
\[
 \frac{iM}{60}\le u\le\frac{(i+1)M}{60},\quad
 \frac{jM}{60}\le w\le\frac{(j+1)M}{60},\quad
 u+w\le M,\qquad i,j\ge0,\ i+j\le59.
\]
There are $60\cdot61/2=1830$ cells per profile. Boundary overlaps are
intentional. On a cell, $v\le M-u_{\rm lo}-w_{\rm lo}$, and
\eqref{app:global:count} with the valid seed $31$ supplies a rational
upper bound for $\sqrt C$. A cell is excluded only if this is strictly
below the profile's $\sqrt C$ lower enclosure.

In each radial row, descend through angular indices until the highest
nonexcluded cell is reached. Its local comparison is evaluated on the
\emph{enlarged closed region}
$0\le u\le u_{\rm hi}$, $0\le w\le w_{\rm hi}$, $u+v+w\le M$.
Consequently one comparison covers all lower angular cells in that
row. This does not assume that coefficients computed on different
regions are monotone. Fully excluded rows retain all their exclusion
inequalities. The complete accounting is
\begin{center}
\begin{tabular}{crrrr}
\toprule
$C$ & Local endpoint & Excluded & Covered by regions & Regions\\
\midrule
$3/4$ & $180$ & $42$ & $1788$ & $60$\\
$11/4$ & $180$ & $1027$ & $803$ & $53$\\
$3$ & $180$ & $908$ & $922$ & $55$\\
$1$ & $182$ & $62$ & $1768$ & $60$\\
\midrule
Total && $2039$ & $5281$ & $228$\\
\bottomrule
\end{tabular}
\end{center}

The odd-profile minimum margins are, respectively,
\[
 \frac{537329859213}{40000000000000}\quad\text{and}\quad
 \frac{3713814057647}{10^{15}}.
\]
The $C=3$ exact margin is a longer rational
stored in its certificate; exact cross multiplication proves it exceeds
$1/10000$ (its decimal value is approximately $0.0001017243329754284$).
For $C=1$ the final lower margin is
\[
 \frac{1420911820069}{500000000000000}>\frac1{10000}.
\]
Thus the uniform margin in \eqref{eq:strict-comparison} is valid.

For $C=3$ we combine the stress-coordinate mean estimate and
finite-normal maximization with the unsigned remainder. For $C=1$ we
use the candidate, localization and normal bounds at endpoint $182$,
the spectral row estimate and the partial-signed remainder bounds.
The four sign weights are $1/4,1/2,3/4,1$. Each row retains all
complete alternatives and selects a valid minimum, with positive
quartic factors and checked cross/same coefficient ordering. There are
$480$ stored scalar tail certificates, each covering $32$ closed
intervals of the degree-six Bernstein polynomial. These bounds include
the mixed terms and the physical radial-weight correction.

The finite verification statement used in
Proposition~\ref{prop:finite-acceptance} is as follows.
For each of the $228$ enlarged regions, independently reconstruct the
candidate bounds, localization interface, physical quadratic coefficients
$c_m,c_r,J,L$, squared normal radius $r_n=\rho_n^2=d_c^2/(2n^2)$, and one complete pair
$A_{\rm rem},B_{\rm rem}$. Acceptance requires all of the following:
\begin{enumerate}
\item The saved closed cells have precisely the declared endpoints,
all signs, vertices and pair types are present, and the count exclusions
and enlarged accepted regions cover all $7320$ budget cells.
\item Every denominator, covariance, phase, rotation and overlap guard
in Sections~\ref{app:global} and~\ref{app:local} holds. In particular
$c_m,c_r>0$, $J,L\ge0$, $0\le r_n\le1/2$,
$A_{\rm rem},B_{\rm rem}\ge0$; the radial floor is positive; and the
required cross/same coefficient orderings hold.
\item The intrinsic and overlap inputs are enclosed on the whole joined
reciprocal domain. In the refined $C=1$ rows the spectral bound uses
$\lambda_+\ge(1+\sqrt3)/4$, which covers $x=0$ as well as the finite
endpoint. Every used signed scalar tail is bounded on its entire disk
by the independently expanded degree-six Bernstein certificate.
\item Select the applicable branch of
\eqref{app:local:normal-max} using these conservative coefficients.
With $g_{\rm normal}$ the negative of its upper enclosure and
${\rm Loss}^+$ an upper enclosure of
\eqref{app:local:loss-formula}, exact rational comparison gives
\begin{equation}\label{app:certificates:accepted-margin}
 g_{\rm normal}-{\rm Loss}^+>\frac1{10000}.
\end{equation}
\end{enumerate}
The arithmetic verifier checks this statement on the fixed witnesses.
Its analytic consequences and propagation are proved in
Section~\ref{app:local:closure}.

\subsection{The worst \texorpdfstring{$C=3$}{C=3} region}
\label{app:certificates:worst}
The minimum over all $55$ accepted $C=3$ regions occurs at row
$(i,j)=(17,23)$. Its original cell has
$u\in[187/1140,33/190]$ and $w\in[253/1140,22/95]$.
The local inequality is certified on the larger set
\begin{equation}\label{app:certificates:worst-domain}
 0\le u\le33/190,\quad 0\le w\le22/95,\quad
 v\ge0,\quad u+v+w\le11/19,
\end{equation}
at the local endpoint $n=180$. Both signs of
$\delta=(-3,1,1,1)/2$ are included in the candidate certificates.
The stress-coordinate mean bound and the physical normal coefficients are
combined with the entire unsigned remainder, including the same-cluster,
radial, mixed and quartic terms. The active normal branch is the interior
branch. The following table gives rational enclosures for this calculation;
the distributed
\path{validation/C3_WORST_ACCEPTANCE.json} gives the exact local inputs,
intermediate coefficients, recorded regional guards and the unabridged
rational margin. Its source-report binding supplies the independently
checked upstream candidate and localization interfaces.
% All decimals are exact rationals.
\begin{center}
\begin{tabular}{lcr}
\toprule
Quantity & Direction & Rational bound\\
\midrule
$c_m$ & $\ge$ & $0.228928718452$\\
$c_r$ & $\ge$ & $0.210023739886$\\
$J$ & $\le$ & $0.815026271354$\\
$L$ & $\le$ & $0.803437886305$\\
$\rho_n^2$ & $\le$ & $0.000132756408$\\
$A_3^{\,u}$ & $\le$ & $0.156826320895$\\
$A_{\rm mixed}$ & $\le$ & $0.018831369020$\\
$A_4^{\,u}$ & $\le$ & $0.014800634464$\\
$A^{\,u}$ & $\le$ & $0.190458324379$\\
$B^{\,u}$ & $\le$ & $0.109690746792$\\
$g_{\rm normal}$ & $\ge$ & $0.205224509235$\\
${\rm Loss}(A^{\,u},B^{\,u})$ & $\le$ & $0.205122784909$\\
$g_{\rm normal}-{\rm Loss}$ & $\ge$ & $0.000101724326$\\
\bottomrule
\end{tabular}
\end{center}
Every terminating decimal in this table denotes its exact rational value.
For the displayed coefficients, exact cross multiplication gives
\[
 J^2\rho_n^2\le(c_m-c_r)^2(1-\rho_n^2),\qquad c_m-c_r>0,
\]
so the interior formula is applicable even after this additional outward rounding.
The displayed $A^{\,u}$ is the sum of its three displayed upper bounds.
The gap and loss rows are recomputed from the displayed coefficients, rather than
copied from a saved final margin. Their difference is strictly greater than
$1/10000$. The unrounded certificate margin is approximately $0.0001017243329754$.
All $55$ regions were compared as exact rationals to identify this minimum.


\subsection{Propagation to all integers in the stated range}
An accepted row fixes its entire region, count floors, fine coefficients,
candidate domain, normal-gap bound, complete remainder alternative,
sign weight and scalar Bernstein bounds. The term-by-term propagation
argument in Section~\ref{app:local:closure}, concluding in
Proposition~\ref{app:local:strict}, shows that the analytic error
majorants retain their endpoint bounds and the permitted normal interval shrinks.
The cross/same majorants remain fixed, so their ordering is preserved.
A later comparison stays within the same accepted scalar disks and
remainder bound. This argument applies throughout the subsequent range
and does not depend on which branch a new evaluation would select.

The two odd profiles apply to all odd $N\ge181$. The $C=3$ profile
applies from $180$, and the $C=1$ profile from $182$; both are therefore
valid for every residue-two integer in the theorem. The exact
divisible case covers $180$ and all later multiples of four. These
branches give every integer $N\ge180$. The local endpoints differ by
profile; the resulting sufficient threshold need not be minimal.

\subsection{Certificates and arithmetic verification}
\label{app:certificates:reproducibility}
The certificates and verifier are publicly available at
\url{https://github.com/motives3/coherence-codimension-two-reproduction}.
This paper refers to the fixed snapshot at
\href{https://github.com/motives3/coherence-codimension-two-reproduction/tree/b7a51dbc324ba260a92c2fafd3642e686cefe91a}{commit \texttt{b7a51dbc324b}}.
All \path{certificate/} and \path{validation/} paths in this paper are
relative to the repository's \path{reproduction/} directory, which is the
reproduction-package root.
The archive \path{certificate/supplement.zip} contains the fourteen JSON
records used by the arithmetic verifier. The machine-readable
\path{certificate/INDEX.json} gives their archive-relative paths, byte
sizes and SHA-256 digests, together with the archive digest. The guide
\path{certificate/ARITHMETIC_README.txt} explains the verifier and its
dependencies. The main adopted records are identified by the following
short names:
\begin{description}
\item[Odd180] The complete subproofs for the two odd profiles in
\texttt{PROFILE\_180.json}, both with certified positive margins.
\item[Stress180] The complete $C=3$ subprofile of
\texttt{STRESS\_SAVED\_NORMAL.json}.
\item[Refined182] \texttt{REFINED\_182\_NORMAL.json}, bound to the immutable
\texttt{FROZEN\_FULL\_182.json} input.
\item[Continuous bounds] The joined root-residual and weighted-overlap
records. The two additional sharp records are
\texttt{SHARP\_FIBERS.json} and \texttt{SHARP\_STATIONARY.json}.
\end{description}
The index supplies the full paths. Some records have an overall failure
label but contain a complete subproof with a certified positive margin.
The verifier checks the adopted subproofs identified above, including
their domains and hypotheses; it does not use the status labels as
acceptance criteria.
The enclosure $\kappa_*$ in Section~\ref{app:local:closure} is stored
in the certificate field \texttt{kappa\_coefficient}; this field
corresponds to the stress coefficient, not the Bloch-Gram coefficient $k_0$.

\begin{samepage}
The verifier consists of four files and uses only the Python standard
library:
\begin{quote}\small
\path{certificate/verify_certificates.py}\\
\path{validation/global_arithmetic.py}\\
\path{validation/continuous_arithmetic.py}\\
\path{validation/local_arithmetic.py}.
\end{quote}
\end{samepage}
\begin{samepage}
Run using Python 3.10 or later:
\begin{quote}\ttfamily\small
python -B certificate/verify\_certificates.py
\end{quote}
\end{samepage}
Run this command from the reproduction-package root. The program locates
its data relative to its own file, reads the ZIP without extraction and
emits JSON to standard output. It creates no files. A failed
guard or arithmetic mismatch exits nonzero. Diagnostic partial stages
are expressly marked incomplete.

The verifier recomputes the global and regional shared budgets,
candidate polar bounds, fixed root intersections, intrinsic residual
bounds, all overlap and projected-fiber intervals, stationary harmonic
rows, and every Bernstein coefficient of the used scalar tails. It then
reconstructs the physical mean/residual curvature, normal maximization,
phase transport and complete local remainder on all $228$ adopted
regions. It does not accept a saved margin as an input to strictness.
The fixed covers and root boxes serve as witnesses, so verification
requires no threshold search, adaptive subdivision or optimization.
The accompanying regression sources exercise the
arithmetic kernels and the rejection of corrupted witnesses.

The historical all-stage report
\path{validation/ARITHMETIC_ACCEPTANCE.json} records the hashes of its
fourteen input records and four verifier files. Their correspondence
with the public distribution is documented in
\path{validation/PRIOR_ACCEPTANCE_BINDING.json};
\path{PRIVACY_TRANSFORM.json} describes the metadata-only redactions.
The current archive and input hashes are listed in \path{certificate/INDEX.json}.

The verifier checks finite rational inequalities. The analytic
identities, geometric implications, infinite-range propagation and
equality classification are proved in the manuscript; they are outside
the verifier's scope. This is a computer-assisted proof, rather than a
proof-assistant formalization.


\section*{Acknowledgments}
\addcontentsline{toc}{section}{Acknowledgments}
OpenAI's Codex was used extensively in the research and preparation of this
manuscript, including use of GPT-6 Astra with the Ultra reasoning setting.
Its use included literature searches and
comparisons with earlier results (Section~\ref{sec:related}); exploration
and development of the four-block construction and its algebraic formulas
(Sections~\ref{sec:construction} and~\ref{app:candidate}); and formulation,
refinement, and critical examination of arguments concerning the matrix
reduction, global localization, local comparison, and equality cases
(Sections~\ref{sec:proof}, \ref{app:global}, and~\ref{app:local}).

Codex was also used in developing the computational work associated
with Section~\ref{app:certificates} and the accompanying reproduction
package. This included implementing and revising code for numerical
exploration, certificate generation, rational-arithmetic verification,
and regression tests.

AI-generated drafts of mathematical arguments and explanatory prose were
used and revised during the preparation of the manuscript. Codex was also used to organize the presentation, revise notation and
cross-references, and edit the English.

The author assumes responsibility for all content, including the
mathematical claims, references, code, and final text.

{\raggedright
\bibliographystyle{plain}
\bibliography{references}
}

\bigskip
\noindent\begin{minipage}{\textwidth}
\small
Chongqing University of Arts and Sciences\\
Chongqing, China\\
Email: \href{mailto:202302474094@stu.cqwu.edu.cn}{\texttt{202302474094@stu.cqwu.edu.cn}}
\end{minipage}
\end{document}
